\documentclass{article}
\usepackage[final]{neurips_2024}
\usepackage[utf8]{inputenc}
\usepackage[T1]{fontenc}
\PassOptionsToPackage{hyphens}{url}
\usepackage{hyperref}
\usepackage{xurl}
\usepackage{booktabs}
\usepackage{amsfonts}
\usepackage{amsmath}
\usepackage{nicefrac}
\usepackage{microtype}
\usepackage{xcolor}
\usepackage{graphicx}
\usepackage{verbatim}
\usepackage{placeins}
\usepackage{needspace}

% Layout overflow mitigations (2026-05-04)
\sloppy
\emergencystretch=3em
\hbadness=10000
\hfuzz=2pt

\title{The Pain Diary and Shadow Discipline:\\A Memory System That Learns from Its Own Incidents}

\author{%
  Luiz Antonio Busnello \\
  Curious Tech Entrepreneur \\
  S\~{a}o Paulo, Brazil \\
  \texttt{lab@generantis.com.br} \\
}

\begin{document}
\raggedbottom

\maketitle

% ---------------------------------------------------------------
% ABSTRACT — sourced from sec_abstract.tex
% ---------------------------------------------------------------
\begin{abstract}
% sec_abstract.tex — Abstract for arXiv submission
% Source: paper-abstract.md (converted 2026-05-04)
% Language: English (paper language)

Production memory systems fail not because retrieval is hard, but because silent
architectural degradation accumulates faster than any embedding model can compensate.
GraphRAG, MemGPT, Mem0, and A-MEM encode structure and recency; none encodes incident
severity, none enforces ranking-change validation. NOX-Supermem's design is
incident-driven: each failure became a schema constraint. Three contributions: (1)
Pain-weighted salience
($\text{salience} = \text{recency} \times \text{pain} \times \text{importance}$,
$\text{pain} \in [0.1, 1.0]$) as a first-class retrieval dimension, calibrated to
PagerDuty P1--P5, not psychometric scales. Ablation
($n=31$, hybrid mode) shows non-significant aggregate effect ($\Delta = +0.0065$, 95\%~CI
$[-0.014, +0.034]$); lift observed in 1/31 queries (Q55, $\Delta = +0.349$); 29/31
unaffected because semantic scores were not tied. The binding constraint
is BM25 recall ceiling: 55 of 60 queries fail
FTS-only regardless of pain calibration; full hybrid retains Recall@10 $= 0.7667$ via
Gemini. (2) Shadow discipline: a seven-day gate before any
ranking change activates, enforced via \texttt{/api/health} as an architectural
constraint, not a convention \cite{kohavi2020trustworthy,chapelle2012interleaved}. (3)
Shared-canonical context: six agents, one corpus, no federation. On 61,257 chunks, hybrid
retrieval (FTS5 + Gemini 3072-d + RRF, $k=60$) achieves
nDCG@10 $= 0.5831 \pm 0.0046$ ($n=60$, 3-run mean, post-cure gold standard v1.1) versus 0.0000 for FTS5 vanilla and
0.1475 for BM25 Pyserini ($n=60$), a $4.0\times$ lift over the strongest lexical baseline.
Edge-type enum coverage improved $14\% \to 56\%$ ($4\times$ gain, $n=100$, 95\%~Wilson
CI $[46$--$66\%]$); self-reported enum coverage rate, not human-validated accuracy. On
LOCOMO (Maharana et al., 2024, $n=100$), FTS5 achieves nDCG@10 $= 0.281$, confirming
lexical difficulty is corpus-dependent; on BEIR TREC-COVID, multilingual-e5-base achieves
nDCG@10 $= 0.8335$ ($n=50$). Open gap: corpus scale $>$100K. Code,
harness, golden set ($n=60$), and incident log:
\url{https://github.com/totobusnello/memoria-nox} (MIT). Operational discipline is at
least as important as embedding sophistication.

\end{abstract}

% ---------------------------------------------------------------
% SECTIONS 1–3 — sourced from sec_1_3.tex
% ---------------------------------------------------------------
% sec_1_3.tex — Sections 1–3 for arXiv submission
% Source: paper-draft-sec1-3.md (converted 2026-05-04)
% Includes §3.8 Pre-registration (not yet in main.tex monolith)
% Language: English (paper language)

% ---------------------------------------------------------------
% SECTION 1 — INTRODUCTION
% ---------------------------------------------------------------
\section{Introduction}

At 22:03 on April 25, 2026, an end-of-day cron job executed \texttt{nox-mem reindex}
without a dry-run flag. Within seconds, 183 entity records lost their \texttt{section},
\texttt{retention\_days}, and \texttt{section\_boost} fields---years of structured
operational context flattened into generic chunks. No error log. No alarm. The database
simply obeyed.

This incident did not expose a retrieval failure. It exposed a discipline failure. The
system had no enforced gate requiring a dry-run preview before a destructive reindex. It
had no mechanism for distinguishing between a trivial documentation update and a production
outage---both were stored with equal weight and retrieved with equal priority. And when six
AI agents (Atlas, Boris, Cipher, Forge, Lex, and Nox) shared a common operational corpus,
a single unchecked operation could silently degrade context for all of them simultaneously.

The incident became the motivating example for this paper.

\subsection{Problem Statement}

LLM agent memory is increasingly recognized as a prerequisite for sustained, coherent
assistance across sessions \cite{lewis2020rag, packer2023memgpt}. Yet the literature on
agent memory systems---from MemGPT's OS-inspired paging \cite{packer2023memgpt} to Mem0's
self-improving extraction \cite{chhikara2025mem0} to A-MEM's Zettelkasten-inspired linking
\cite{xu2025amem}---converges on a narrow set of concerns: retrieval accuracy on benchmark
corpora, latency at scale, and memory capacity. What these systems do not address is the
operational discipline required to keep a production memory system trustworthy over months
of continuous use.

Three specific gaps drive this work:

\textbf{Silent regression risk.} Ranking changes---new scoring formulas, boosting
coefficients, retrieval layer modifications---can degrade retrieval quality without any
observable error signal. Production memory systems have no established methodology for
validating such changes before they affect live queries. The April 25 incident is one
instance of a broader pattern: operations that silently alter system behavior with no guard
and no rollback.

\textbf{Incident severity as an invisible dimension.} Existing memory systems model recency
and semantic relevance. None, to our knowledge, model the operational cost of the
experience being stored. A lesson learned from a production outage and a routine
implementation note may share similar timestamps and embedding distances; they do not share
similar importance to future retrieval. This asymmetry is unrepresented in current
retrieval architectures.

\textbf{Multi-agent context silos.} When multiple agents operate over isolated memory
stores, intelligence accumulated by one agent is inaccessible to others. Solutions that
partition by \texttt{user\_id} \cite{chhikara2025mem0} or maintain per-agent state
\cite{packer2023memgpt} impose federation overhead or simply accept the silo. For trusted
multi-agent deployments---where agents serve a single operator with a shared context---this
isolation is an architectural choice, not a requirement.

\subsection{Thesis}

\textbf{Agent memory is not a retrieval problem---it is an operational discipline problem.}

Better embeddings, larger corpora, and faster fusion layers address the retrieval surface.
They do not address the question of whether the system can be trusted to evolve without
silent regressions, whether it surfaces the most operationally important context rather
than merely the most recent, or whether the intelligence accumulated by one agent is
available to the others who need it.

\subsection{Contributions}

We present \textbf{nox-mem}, a production memory system designed for multi-agent
operational use and built over approximately three months of continuous deployment.\footnote{Repository:
\url{https://github.com/totobusnello/memoria-nox}; submission tar at
\texttt{paper/publication/}; license MIT.} The system currently indexes 61,257 chunks
across a shared canonical corpus, maintains 99.97\% vector coverage, and operates at p95
latency below one second on commodity VPS hardware (8-core x86\_64) over the production
query mix. Three primary contributions distinguish this work:

\begin{enumerate}

\item \textbf{Pain-weighted salience as a secondary ranking modulator} (\S3.5). We
\textbf{propose} a retrieval scoring signal that explicitly models incident severity:
$\texttt{salience} = \texttt{recency} \times \texttt{pain} \times \texttt{importance}$,
where $\texttt{pain} \in [0.1, 1.0]$ encodes the operational cost of the recorded
experience---from trivial notes (0.1) to production outages (1.0). To our knowledge, this
is the first documented retrieval signal in the LLM agent memory systems literature that
treats incident severity as a first-class dimension; severity-aware prioritization itself
has well-established roots in observability and incident-response practice
(e.g., PagerDuty severity tiers, SIEM alert weighting), but has not been instrumented as
a retrieval ranking signal in the memory-systems literature surveyed in \S2.
Empirical ablation (\S5.4, E10, $n=31$
post-incident queries, hybrid mode) shows directional but non-significant aggregate effect
($\Delta = +0.0065$, 95\% CI $[-0.014, +0.034]$). The dimension provides meaningful lift
in the narrow regime where high-pain and low-pain chunks receive near-identical Gemini
semantic scores (Q55 case study, $\Delta = +0.349$). The transferable contribution is the
methodology of treating incident severity as typed schema input---a design pattern adoptable
independently of any retrieval stack---rather than a demonstrated corpus-wide retrieval
improvement.

\item \textbf{Enforced shadow discipline} (\S3.6). We describe a methodology for
validating ranking changes against a shadow telemetry baseline before deployment,
implemented as an architectural constraint rather than a documented suggestion. The
mechanism is codified via environment variable gating
(\texttt{NOX\_SALIENCE\_MODE=shadow|active}), enforced by a cron-checked health endpoint,
with a mandatory minimum observation period of seven days. During the salience validation
run (Fase~1.7b-b), this produced 191 promotion candidates, 16,608 review cases, and
45,743 archive recommendations before any activation. This is, to our knowledge, the first
memory system paper to codify such a discipline as an architectural constraint rather than
a suggested practice.

\item \textbf{Shared-canonical multi-agent design} (\S3.7). Six agents share a single
\texttt{chunks} table distinguished by \texttt{source\_file} prefix rather than separate
stores. Cross-agent retrieval requires no synchronization because partitioning never
occurred. When Forge records a decision about infrastructure, Atlas retrieves it directly
on the next relevant query---not because data was replicated, but because the design
assumed shared context from the start.

\end{enumerate}

Secondary contributions include: a reproducible evaluation harness with nDCG@10, MRR,
Recall@10, and Precision@5 over 50 curated golden queries; approximately three months of operational
evidence across five infrastructure upgrades without data loss; and comparative analysis
against strong retrieval baselines (see \S5).

\subsection{Paper Organization}

Section~2 surveys related work across knowledge graph--augmented retrieval, agent memory
systems, and production-oriented frameworks, identifying the specific gaps this work
addresses. Section~3 describes the system architecture across storage, retrieval, knowledge
graph, and operational discipline layers. Section~4 details the evaluation methodology,
including the golden query harness, baseline selection, and three-corpus design. Section~5
presents experimental results, including hybrid pipeline ablations, baseline comparisons,
pain dimension validation, and the shadow discipline case study. Section~6 discusses
deployment scope, limitations, and future work. Section~7 concludes.

% ---------------------------------------------------------------
% SECTION 2 — RELATED WORK
% ---------------------------------------------------------------
\section{Related Work}

Agent memory research has accelerated rapidly since the introduction of
retrieval-augmented generation \cite{lewis2020rag}, with recent work spanning knowledge
graph--augmented retrieval, autonomous memory management, and production deployment
frameworks. We survey the most directly relevant systems and identify the specific
dimensions where this work diverges.

\subsection{Knowledge Graph--Augmented Retrieval}

\textbf{GraphRAG} \cite{edge2024graphrag} introduced a pipeline that uses LLM extraction
to build entity-relation graphs over document corpora, applies community detection (Leiden
algorithm) to identify thematic clusters, and generates multi-level summaries for
query-focused retrieval. The approach substantially improves global query coverage---queries
requiring synthesis across many documents---compared to naive RAG. nox-mem adopts the
LLM-extraction paradigm for knowledge graph construction but diverges in objective:
GraphRAG optimizes for global summarization via community traversal, while nox-mem targets
entity-centric local retrieval for operational contexts, with closed-enum edge typing
designed to support downstream blast-radius analysis. GraphRAG does not model retrieval as
an operational concern and has no analogue to the discipline mechanisms described in
\S3.5--3.6.

\textbf{HiRAG} \cite{huang2025hirag} extends the KG-augmented paradigm with hierarchical
knowledge representation and iterative reasoning chains, demonstrating state-of-the-art
performance on multiple QA benchmarks. The hierarchical approach---traversing from coarse to
fine-grained graph levels---optimizes for accuracy on complex reasoning tasks. nox-mem
instead maintains a flat three-layer fusion pipeline (\S3.4), accepting the
accuracy-latency trade-off explicitly: p95 below one second at 61K chunks serves
operational use cases where query latency is a usability constraint, not just a benchmark
metric. HiRAG does not address multi-agent deployment, evaluation harnesses for operational
corpora, or retrieval discipline.

\textbf{HippoRAG} \cite{guo2024hipporag} draws on neurobiological memory models to improve
associative retrieval via hippocampal-inspired graph structures. The system demonstrates
that non-uniform memory consolidation---weighting some associations more strongly than
others based on structural graph properties---improves retrieval over flat indexing. nox-mem
introduces a conceptually related asymmetry through the pain dimension (\S3.5), though the
design rationale is grounded in operational incident management practice rather than
neuroscience. HippoRAG does not model incident severity as a retrieval signal.

\subsection{Agent Memory Systems}

\textbf{MemGPT} \cite{packer2023memgpt} proposed treating the LLM itself as an operating
system, with the agent autonomously managing its own context through function calls that
page between main context and external storage. The architecture gives agents direct
control over their memory---a compelling design for single-agent autonomy. nox-mem adopts
the opposite philosophy: memory is managed by an external pipeline (file watcher, ingest
router, nightly KG extraction), and agents are consumers of a shared canonical store rather
than managers of isolated state. This separation enables cross-agent retrieval as a
structural property while imposing the discipline that agents cannot corrupt their own
context.

\textbf{Mem0} \cite{chhikara2025mem0} introduced a production-focused memory layer with
LLM-driven memory extraction, deduplication, and temporal updates. Evaluated on the LOCOMO
benchmark, Mem0 demonstrates +26\% accuracy improvement over baselines on long-context
conversation scenarios. nox-mem shares the production emphasis but diverges in two
important respects. First, Mem0 relies on autonomous LLM editing decisions to update and
merge memories; nox-mem requires shadow-period validation before any retrieval-affecting
change propagates to production, specifically to mitigate silent regression risk. Second,
Mem0's multi-user design partitions memory by \texttt{user\_id}, treating agents as
user-addressable endpoints; nox-mem's shared canonical design treats agents as roles within
a single trusted operational context.

\textbf{A-MEM} \cite{xu2025amem} introduced agentic memory with Zettelkasten-inspired
structure: notes are automatically tagged, interlinked, and contextualized by an LLM
controller, enabling dynamic memory evolution. The approach handles unstructured
conversation well but does not address operational corpora with explicit schema
requirements, retention policies, or incident-severity differentiation. A-MEM has no
evaluation harness independent of LLM-generated quality judgments, and does not address
multi-agent deployments or ranking change discipline.

\subsection{Production-Oriented Frameworks}

\textbf{Cognee} \cite{topoteretes2024cognee} provides an open-source AI memory framework
implementing an Extract-Cognify-Load (ECL) pipeline with native knowledge graph support
and hybrid retrieval. Of the systems surveyed, Cognee is architecturally closest to
nox-mem: both maintain native KGs, both support hybrid retrieval, and both target
production deployment scenarios. Cognee achieves 3/5 on our architectural parity dimensions
(Table~\ref{tab:comparison}). The key differentiators are shadow discipline---Cognee provides
no enforced validation gate for retrieval-affecting changes---and the evaluation harness:
Cognee lacks a reproducible evaluation methodology with standard IR metrics over a fixed
golden set. nox-mem's eval harness (\S4) enables the kind of regression detection that the
shadow discipline gate is designed to act on.

\textbf{LangChain Memory} provides key-value and buffer memory primitives widely used in
production applications. The design prioritizes API simplicity and framework compatibility
over retrieval quality; it does not support hybrid retrieval, knowledge graph construction,
or evaluation. It serves as a practical lower bound in comparative analysis.

\subsection{Foundational Retrieval Methods}

\textbf{Retrieval-Augmented Generation} \cite{lewis2020rag} established the foundational
architecture combining parametric LLM knowledge with non-parametric document retrieval,
enabling factual question answering over dynamic corpora. nox-mem extends this foundation
to persistent, evolving operational corpora with explicit retention, pain weighting, and
multi-agent routing.

\textbf{Reciprocal Rank Fusion} \cite{cormack2009rrf} provides the rank-combination
mechanism at the core of nox-mem's hybrid retrieval layer. Cormack et al.\ demonstrated
that RRF consistently outperforms individual rank learning methods and Condorcet-based
fusion without requiring system-specific tuning. We adopt RRF with $k=60$ as the fusion
parameter, consistent with the paper's recommendations, as the final combination step over
BM25 and dense retrieval ranked lists.

\subsection{Positioning: Where nox-mem Fits}
\label{sec:positioning}

Table~\ref{tab:comparison} maps seven architectural and operational axes across the
surveyed systems. The table is intended to illustrate where the systems differ
structurally, not to function as a competition score. We include two dimensions where
nox-mem does not yet have coverage---corpus scale above 100K chunks and third-party
benchmark evaluation---because honest comparison requires acknowledging these gaps
explicitly. Limitations arising from these gaps are addressed in \S6.3; planned mitigations
are in \S6.5.

\begin{table}[!ht]
  \caption{Architectural comparison across memory systems. Axes reflect structural design
    choices, not performance rankings. \checkmark\ = implemented; $\approx$ = partial or
    indirect; $\times$ = absent.}
  \label{tab:comparison}
  \centering
  \small
  \resizebox{\textwidth}{!}{%
\begin{tabular}{llllll}
    \toprule
    System & KG Native & Hybrid Retr. & Eval Harness & Multi-Agent & Shadow Disc. \\
    \midrule
    \textbf{nox-mem (ours)} & closed-enum, 7 & FTS5+Gemini+RRF & nDCG/MRR/Recall & shared canonical & enforced $\geq$7d \\
    GraphRAG \cite{edge2024graphrag} & community det. & via KG queries & none & none & none \\
    MemGPT \cite{packer2023memgpt} & none & embedding-first & none & per-agent state & none \\
    Mem0 \cite{chhikara2025mem0} & optional (v2) & vector-only & LOCOMO only & user\_id part. & none \\
    A-MEM \cite{xu2025amem} & Zettelkasten & semantic-first & none & none & none \\
    HiRAG \cite{huang2025hirag} & hierarchical & multi-level & task-specific & none & none \\
    Cognee \cite{topoteretes2024cognee} & ECL pipeline & hybrid & ad-hoc & optional & none \\
    LangChain Memory & none & key-value & none & session\_id & none \\
    \bottomrule
  \end{tabular}}
  \begin{minipage}{\linewidth}
    \vspace{4pt}\footnotesize
    Five architectural axes are displayed; two additional axes---\emph{corpus scale
    \textgreater100K chunks} and \emph{3rd-party benchmark evaluation} (e.g., LOCOMO,
    LongMemEval, BEIR)---are summarized in the prose below. nox-mem currently covers
    5 of 7 axes ($\times$ on both omitted axes; see \S5.2 for partial BEIR/LOCOMO
    progress). Mem0 covers axis 7 via LOCOMO; no surveyed system has published a
    \textgreater100K-chunk evaluation. This is a comparison across architectural
    \emph{axes of difference}, not a competition score.
  \end{minipage}
\end{table}

We compare across seven axes. nox-mem covers 5/7: knowledge graph, hybrid retrieval, eval
harness, multi-agent design, and enforced shadow discipline. It does \textbf{not} yet cover
corpus scale above 100K chunks (current corpus: 61K) or evaluation against third-party
benchmarks such as LOCOMO, LongMemEval, or BEIR---the evaluation harness uses an internally
curated golden set. These are explicit limitations, not omissions, and are documented in
\S6.3 and \S6.5. The two axes with zero coverage across all surveyed systems---pain-weighted
salience and shadow discipline---are the primary novelty claims of this paper. Shadow
validation as a technique has well-established roots in production search evaluation
\cite{chapelle2012interleaved} and online controlled experiment methodology
\cite{kohavi2020trustworthy}. Our contribution is the instantiation of this discipline as
an enforced architectural constraint in an LLM agent memory system, rather than as a
testing methodology applied to web-scale traffic populations. Section~3 describes the
nox-mem architecture across all five covered dimensions; Sections~4--5 provide empirical
grounding.

% ---------------------------------------------------------------
% SECTION 3 — SYSTEM ARCHITECTURE
% ---------------------------------------------------------------
\section{System Architecture}

nox-mem is a production memory system built on SQLite with three primary subsystems---storage,
retrieval, and knowledge graph---plus an operational discipline layer that governs how the
system evolves. The implementation is in TypeScript using \texttt{better-sqlite3} for
synchronous database access, \texttt{sqlite-vec} for approximate nearest-neighbor search,
and Gemini for embedding generation and LLM-based KG extraction. The system has operated
continuously since March 2026 with 12 schema migrations and no irrecoverable data loss (one
metadata-loss event on 2026-04-25 was recovered via re-ingestion; see \S6.2).

\subsection{System Overview}

\begin{figure}[!ht]
  \centering
  \includegraphics[width=\linewidth]{figures/figure1-system-overview.pdf}
  \caption{System overview. Six personas read and write through a unified interface layer
    (CLI / MCP / HTTP) into a single canonical SQLite corpus (no per-agent silos).
    Retrieval runs as a three-layer hybrid (BM25 $\to$ 3072-d Gemini cosine $\to$ RRF,
    $k=60$). The shadow-mode pipeline and append-only \texttt{ops\_audit} are cross-cutting
    governance: every ranking change ships shadow-first; every destructive op is wrapped by
    \texttt{withOpAudit()} snapshot.}
  \label{fig:overview}
\end{figure}
\FloatBarrier

The entry point for all CLI operations is \texttt{dist/index.js} (the compiled package
binary). An HTTP API on port 18802 exposes search, knowledge graph, health, and operational
endpoints for agent integration.

\medskip
\begin{verbatim}
Source files (.md / .docx / .pptx / .pdf / entity files)
    |
    v
ingest-router (inotifywait watcher + manual CLI)
    |
    +------> chunks table (canonical, 61K rows)
                |
                +------> chunks_fts (FTS5 BM25 index)
                +------> vec_chunks (sqlite-vec, 3072d embeddings)
                +------> kg_entities + kg_relations (LLM-extracted, nightly)
                +------> search_telemetry (shadow-mode + eval harness)
                |
    query --> search() pipeline:
                Layer 1: FTS5 BM25 candidates
                Layer 2: Gemini semantic candidates
                Layer 3: RRF fusion (k=60) + salience boost + section_boost
                --> ranked results (p95 < 1s @ 61K chunks)
\end{verbatim}
\medskip

\subsection{Storage Layer}

The canonical store is a single SQLite database accessed via \texttt{better-sqlite3} for
synchronous, blocking reads---a deliberate choice that simplifies the concurrency model for
a single-node deployment. Three coordinated tables implement the storage layer.

\texttt{chunks} is the central table, containing content text, metadata
(\texttt{source\_file}, \texttt{chunk\_type}, \texttt{created\_at}, \texttt{last\_seen}),
and operational fields (\texttt{pain} REAL DEFAULT 0.2, \texttt{importance} REAL,
\texttt{section} TEXT, \texttt{section\_boost} REAL, \texttt{retention\_days}). The schema
has evolved through 12 versions (v1 through v12) since March 2026; all migrations have been
applied to the live database without downtime or data loss. Key schema additions include
\texttt{retention\_days} (v8, typed retention by chunk category), \texttt{pain} (v9,
severity signal), and \texttt{section} with \texttt{section\_boost} (v10, structural
weighting for entity file sections).

\texttt{chunks\_fts} is an FTS5 virtual table providing BM25-ranked full-text search over
chunk content. FTS5 vanilla AND-mode behavior---requiring all query terms to
match---means that natural language queries frequently return zero results when run against
the lexical index alone. This structural characteristic of FTS5 makes hybrid retrieval a
necessity rather than an optimization, a finding confirmed in evaluation (\S5.1).

\texttt{vec\_chunks} stores 3072-dimensional Gemini embeddings
(\texttt{gemini-embedding-001}) via the \texttt{sqlite-vec} extension, enabling
approximate cosine similarity search without external vector store infrastructure. Current
vector coverage is 99.97\% across 61,257 chunks.

Retention policy is encoded directly in the schema via \texttt{retention\_days}: feedback
and person records are set to \texttt{NULL} (never decay); decisions and projects have
365-day retention; team context 120 days; daily notes 90 days; pending items 30 days. This
ensures that the system's own operating rules cannot be silently evicted by a generic decay
policy.

\subsection{Retrieval Layer}

The retrieval pipeline operates in three sequential layers, with results from layers one and
two merged via Reciprocal Rank Fusion \cite{cormack2009rrf}:

\textbf{Layer 1---Lexical retrieval.} FTS5 BM25 search expands the query against the
\texttt{chunks\_fts} index, returning up to $k_1$ ranked candidates. Results are
deduplicated by chunk ID before passing to fusion.

\textbf{Layer 2---Semantic retrieval.} The query is embedded using
\texttt{gemini-embedding-001} (3072-d), and cosine similarity search is performed against
\texttt{vec\_chunks} via \texttt{sqlite-vec}, returning up to $k_2$ ranked candidates.

\textbf{Layer 3---RRF fusion and boosting.} Candidates from layers one and two are merged
via Reciprocal Rank Fusion \cite{cormack2009rrf} with $k=60$:

\[
  s_{\text{rrf}}(d) = \frac{1}{k + r_{\text{fts}}(d)} + \frac{1}{k + r_{\text{sem}}(d)}
\]

where $r_{\text{fts}}(d)$ and $r_{\text{sem}}(d)$ are the BM25 and cosine-similarity ranks
of chunk $d$ respectively.

The fused score is then combined with two boost signals via log-additive aggregation:

\[
  \log s_{\text{final}}(d) = \log s_{\text{rrf}}(d) + \log s_{\text{salience}}(d)
    + \log s_{\text{section}}(d)
\]

where $s_{\text{salience}}(d) = \texttt{recency}(d) \times \texttt{pain}(d) \times
\texttt{importance}(d) \in [0, 2]$ (\S3.5), and $s_{\text{section}}(d) \in
\{2.0,\,1.5,\,0.8,\,1.0\}$ for compiled, frontmatter, timeline, and unstructured (NULL)
section types respectively.

Equivalently in linear space,
$s_{\text{final}}(d) = s_{\text{rrf}}(d) \cdot s_{\text{salience}}(d) \cdot s_{\text{section}}(d)$.
The two formulations are mathematically identical; the log-additive form is operationally
preferred because it allows each boost's contribution to be recorded and monitored
independently in the \texttt{search\_telemetry} log without re-multiplication.

Top-$k$ chunks are returned in descending order of $s_{\text{final}}$.

The pipeline achieves p95 search latency below one second across the 61K chunk corpus,
measured in production over the observation period. The \texttt{--no-hybrid} flag disables
semantic retrieval for latency-constrained contexts where lexical precision is acceptable.

\subsection{Knowledge Graph Layer}

An LLM-extracted knowledge graph augments chunk retrieval with structured relationship
data. Entities and relations are extracted incrementally via nightly batch runs using
\texttt{gemini-2.5-flash} and stored in \texttt{kg\_entities} ($\approx$402 entities) and
\texttt{kg\_relations} ($\approx$544 relations).

Entity types follow a closed taxonomy of 10 categories: person, project, agent, tool,
concept, organization, technology, document, decision, and metric. Relation edges use
free-form \texttt{relation\_type} text combined with a closed-enum
\texttt{relation\_reason} field with seven canonical values: \texttt{depends\_on},
\texttt{derived\_from}, \texttt{opposes}, \texttt{extends}, \texttt{replaces},
\texttt{mentions}, and \texttt{unknown}. The closed enum enables downstream
tooling---\texttt{impact~\textless entity\textgreater} computes one-hop blast radius with
reason-priority weighting; \texttt{detect-changes} maps git diffs to affected KG
entities---without requiring LLM inference at query time.

Edge type enum coverage required explicit engineering. Initial extraction with an optional
enum field and a ``use unknown if unsure'' prompt instruction produced 86\% \texttt{unknown}
labels across n=100 sampled relations. A three-path fix---a revised prompt with explicit
examples for each non-unknown category; a code-side defensive normalization map
(\texttt{RELATION\_TYPE\_TO\_REASON}, 24 input aliases covering PT-BR and EN free-text
variants that collapse to the 7 canonical reasons); and a post-extraction validation
pass---improved the enum coverage rate from 14\% to 56\% ($4\times$ improvement),
equivalently reducing the \texttt{unknown} rate from 86\% to 44\% (n=100 sample). The
enum coverage rate measures the proportion of sampled relations where the LLM committed to
one of the six non-unknown typed values rather than falling back to \texttt{unknown}; it is
a self-reported coverage rate, not a classification accuracy against a human-annotated
ground truth (see \S6.3). This experience reinforces a general principle: LLM-extracted
structured fields with optional enums and uncertainty-deferring prompts will systematically
underperform; defensive code-side normalization is required.

\subsection{Pain-Weighted Salience}

The salience formula is:

\begin{equation}
  \texttt{salience}(\mathit{chunk}) =
    \texttt{recency}(\mathit{chunk})
    \times \texttt{pain}(\mathit{chunk})
    \times \texttt{importance}(\mathit{chunk})
  \label{eq:salience}
\end{equation}

\noindent where:
\begin{itemize}
  \item $\texttt{recency} \in [0,\,1]$ is an exponential decay function over the
    \texttt{last\_seen} timestamp, parameterized by the chunk's \texttt{retention\_days}
    value;
  \item $\texttt{pain} \in [0.1,\,1.0]$ encodes the operational severity of the recorded
    experience, annotated via \texttt{pain: 0.X} markers in entity files. Values are
    assigned on a five-point semantic scale: 0.1 (trivial note), 0.3 (minor friction),
    0.5 (significant issue), 0.7 (major incident), 1.0 (production outage with data loss
    or downtime);
  \item $\texttt{importance} \in [0,\,1]$ is derived from \texttt{mention\_count} across
    the corpus and entity-type priors.
\end{itemize}

The formula is multiplicative: a high-pain chunk with low recency can still score above a
low-pain chunk with high recency when the severity differential is large. This reflects an
engineering choice motivated by operational practice rather than a biological or
psychometric claim. The five-point scale follows the structure of established incident
severity taxonomies \cite{pagerduty2023severity,beyer2016site}: just as incident management
systems assign higher dispatch priority to a production outage than a documentation note
regardless of their timestamps, pain weighting ensures that incident-derived memory
dominates retrieval over routine content even when recency is comparable.

The salience signal is exposed in shadow mode at \texttt{/api/health.salience} without
affecting production ranking, enabling comparison against the baseline before activation
(\S3.6). Ablation results comparing retrieval quality with real versus uniform pain values
on $n=31$ post-incident queries are reported in \S5.4 (E10, 2026-05-04); the aggregate
result is directional but not statistically significant, with Q55 as the primary positive
case study ($\Delta = +0.349$).

\FloatBarrier
\subsection{Enforced Shadow Discipline}

\begin{figure}[!ht]
  \centering
  \includegraphics[width=0.9\linewidth]{figures/figure3-shadow-state-machine.pdf}
  \caption{Shadow discipline state machine. Every ranking-affecting change traverses this
    machine. Activation requires three independent signals: at least seven days of shadow
    telemetry, a healthy distribution (no degenerate clustering), and explicit human
    approval. Regression at any point reverts via a single environment-variable
    flip---zero schema rollback required.}
  \label{fig:shadow}
\end{figure}
\FloatBarrier

The shadow discipline mechanism enforces a separation between observing a retrieval change
and deploying it. Any modification that affects ranking---scoring formula changes, boost
adjustments, new retrieval layers---must run in shadow mode for a minimum of seven days
before activation. This is an architectural constraint, not a documented recommendation.
Shadow validation has well-established roots in production search and online controlled
experiments \cite{kohavi2020trustworthy,chapelle2012interleaved}. Our contribution is the
application of shadow discipline specifically to LLM agent memory systems, with
architectural enforcement via cron and health endpoints.

Implementation relies on three components: (1)~an environment variable gate
(\texttt{NOX\_SALIENCE\_MODE=shadow|active}) that controls whether the salience score
affects production ranking or is only written to \texttt{search\_telemetry}; (2)~a health
endpoint (\texttt{/api/health.opsAudit}) that cron checks every 15 minutes and alerts via
Discord if the shadow observation period has not been met; (3)~an append-only audit log
(\texttt{ops\_audit} table) with database triggers that block DELETE and UPDATE on rows
with terminal status, ensuring the shadow-period record cannot be retroactively modified.

The distinction from A/B testing is important: A/B testing requires live traffic over a
production population. Shadow discipline for operational/personal corpora has no such
traffic---queries are infrequent, irregular, and non-i.i.d. The shadow mechanism instead
collects telemetry on what the proposed ranking \emph{would have done} on each live query,
accumulating promotion, review, and archive candidate distributions across the observation
period. The distribution comparison then drives the activation decision.

The salience validation run (Fase~1.7b-b) ran for seven days before activation, producing
191 promotion candidates, 16,608 review cases, and 45,743 archive recommendations. The
distribution was reviewed, found consistent with design intent, and activation proceeded.
Counterfactual analysis suggests that the April~25 reindex incident---the motivating example
in \S1---would have been detected during a shadow period: the reindex operation altered
chunk distributions in ways that would have surfaced as anomalous in telemetry within hours.

\subsection{Shared-Canonical Multi-Agent Design}

Six agents (Atlas, Boris, Cipher, Forge, Lex, and Nox) share a single \texttt{chunks}
table. Agent-specific content is distinguished by \texttt{source\_file} path prefix
(\texttt{agents/<name>/...}) rather than separate tables, schemas, or databases.
Cross-agent queries use \texttt{cross-search} and \texttt{cross-kg} CLI commands, which
apply additional \texttt{source\_file} filters to retrieve context generated by a specific
agent.

The design assumes a single trusted operator. It is explicitly not suitable for
multi-tenant SaaS deployments, where user isolation is a security requirement rather than
an engineering choice. Within the trusted-operator constraint, shared-canonical design
provides cross-agent intelligence as a structural property: when Forge ingests a decision
about an infrastructure change, Atlas retrieves it on the next architecturally relevant
query without any synchronization step.

The operational tooling layer---\texttt{detect-changes}, \texttt{impact},
\texttt{api-impact}, \texttt{consolidate-merge}, \texttt{kg-reclassify}---extends the
system to support multi-agent workflows beyond retrieval, enabling agents to query the KG
for blast-radius analysis before executing changes, and to surface merge candidates across
agent-sourced entities.

The architecture in \S3 provides the substrate for the empirical evaluation in \S4--5.
Section~4 describes the evaluation methodology, including the golden query harness,
baseline systems, and three-corpus protocol designed to address the single-corpus and
single-curator limitations of operational memory benchmarking.

\FloatBarrier
\subsection{Pre-registration and Reproducibility Commitments}

A persistent concern in machine learning evaluation is post-hoc query set construction:
a curated golden set assembled with knowledge of retrieval results can unintentionally
favor the system under evaluation, inflating reported metrics without any deliberate
manipulation \cite{munafosystematic2017}. We address this concern explicitly by anchoring
the experimental design to a pre-registration artifact that predates all result collection.

\textbf{Pre-registration document.} The evaluation protocol---including query categories,
metrics, baseline selection criteria, and success thresholds for each
hypothesis---was committed to \texttt{paper/publication/03-experiments-needed.md} before
any baseline or ablation experiment was executed. This document is tracked in the public
repository (\url{https://github.com/totobusnello/memoria-nox}) and its commit history is
immutable. The git commit hash and ISO-8601 author timestamp for that file should be
verified at submission via:

\medskip
\begin{verbatim}
git log --follow -1 --format="%H  %aI" \
    paper/publication/03-experiments-needed.md
\end{verbatim}
\medskip

\textbf{Golden query set freeze.} The 60-query golden set (\texttt{eval/golden-queries.jsonl},
comprising the 50-query main set R01b and the 10-query held-out subset R01c) is included
in the public repository and the submission tar. The verifiable identifiers are:

\begin{itemize}
  \item \textbf{Git commit hash:} \texttt{f75d1864b581d0a0933704e997669bbd41aa507f}
  \item \textbf{Author timestamp (ISO-8601, BRT):} \texttt{2026-05-04T13:38:01-03:00}
  \item \textbf{Freeze mtime (VPS, UTC):} \texttt{2026-05-04T09:45Z}
  \item \textbf{SHA-256:} \texttt{9bff8ee7b9056eff6a1af22305cae762aa4b98e682578faffb4c22cdd0a2cd7d}
  \item \textbf{File size:} 8990 bytes, 60 lines (one query per line, JSONL)
\end{itemize}

The SHA-256 hash allows any third party to verify byte-for-byte that the included
\texttt{golden-queries.jsonl} matches the artifact used for all reported evaluations.
The golden query set was frozen \emph{before} any BM25 (Pyserini), multilingual-e5-base,
BEIR TREC-COVID, LOCOMO, or pain-ablation (E10) experiment was executed; all results in
\S5 were produced against this exact file.

\textbf{Held-out subset disclosure.} The 10 queries designated R01c (held-out for baseline
robustness validation) were locked as a named subset within the golden set before the final
tuning sprint that fixed the RRF $k=60$ parameter and the salience boost coefficients.
They were not constructed after tuning; they are a partition of the original 60-query set,
identified by their position index (Q51--Q60) in the JSONL file. R01c is evaluated at most
once per major system revision and its results are reported separately from the 50-query
main set to prevent iterative query refinement bias (\S4.1).

\textbf{Registered reports context.} While this paper does not constitute a formal
registered report in the sense of \cite{munafosystematic2017}---no journal or conference
pre-acceptance was obtained prior to experimentation---the \texttt{03-experiments-needed.md}
document serves an equivalent function within the project's operational discipline: it
records pre-registered hypotheses with explicit success thresholds (e.g.,
$\Delta$~nDCG@10 $\geq 0.05$ for E10) that are reported honestly regardless of outcome,
as evidenced by the DIRECTIONAL, NOT SIGNIFICANT result for the E10 pain ablation (\S5.4).
The pre-registration document, the golden query set, and all evaluation scripts are
included in the public repository under MIT license, permitting full independent
verification.


% [ORIGINAL MONOLITH PRESERVED BELOW FOR REFERENCE — remove before submission]
% ---------------------------------------------------------------
% SECTION 1 — INTRODUCTION (MONOLITH)
% ---------------------------------------------------------------
\iffalse
\section{Introduction}

At 22:03 on April 25, 2026, an end-of-day cron job executed \texttt{nox-mem reindex}
without a dry-run flag. Within seconds, 183 entity records lost their \texttt{section},
\texttt{retention\_days}, and \texttt{section\_boost} fields---years of structured
operational context flattened into generic chunks. No error log. No alarm. The database
simply obeyed.

This incident did not expose a retrieval failure. It exposed a discipline failure. The
system had no enforced gate requiring a dry-run preview before a destructive reindex. It
had no mechanism for distinguishing between a trivial documentation update and a production
outage---both were stored with equal weight and retrieved with equal priority. And when six
AI agents (Atlas, Boris, Cipher, Forge, Lex, and Nox) shared a common operational corpus,
a single unchecked operation could silently degrade context for all of them simultaneously.

The incident became the motivating example for this paper.

\subsection{Problem Statement}

LLM agent memory is increasingly recognized as a prerequisite for sustained, coherent
assistance across sessions \cite{lewis2020rag, packer2023memgpt}. Yet the literature on
agent memory systems---from MemGPT's OS-inspired paging \cite{packer2023memgpt} to Mem0's
self-improving extraction \cite{chhikara2025mem0} to A-MEM's Zettelkasten-inspired linking
\cite{xu2025amem}---converges on a narrow set of concerns: retrieval accuracy on benchmark
corpora, latency at scale, and memory capacity. What these systems do not address is the
operational discipline required to keep a production memory system trustworthy over months
of continuous use.

Three specific gaps drive this work:

\textbf{Silent regression risk.} Ranking changes---new scoring formulas, boosting
coefficients, retrieval layer modifications---can degrade retrieval quality without any
observable error signal. Production memory systems have no established methodology for
validating such changes before they affect live queries. The April 25 incident is one
instance of a broader pattern: operations that silently alter system behavior with no guard
and no rollback.

\textbf{Incident severity as an invisible dimension.} Existing memory systems model recency
and semantic relevance. None, to our knowledge, model the operational cost of the
experience being stored. A lesson learned from a production outage and a routine
implementation note may share similar timestamps and embedding distances; they do not share
similar importance to future retrieval. This asymmetry is unrepresented in current
retrieval architectures.

\textbf{Multi-agent context silos.} When multiple agents operate over isolated memory
stores, intelligence accumulated by one agent is inaccessible to others. Solutions that
partition by \texttt{user\_id} \cite{chhikara2025mem0} or maintain per-agent state
\cite{packer2023memgpt} impose federation overhead or simply accept the silo. For trusted
multi-agent deployments---where agents serve a single operator with a shared context---this
isolation is an architectural choice, not a requirement.

\subsection{Thesis}

\textbf{Agent memory is not a retrieval problem---it is an operational discipline problem.}

Better embeddings, larger corpora, and faster fusion layers address the retrieval surface.
They do not address the question of whether the system can be trusted to evolve without
silent regressions, whether it surfaces the most operationally important context rather
than merely the most recent, or whether the intelligence accumulated by one agent is
available to the others who need it.

\subsection{Contributions}

We present \textbf{nox-mem}, a production memory system designed for multi-agent
operational use and built over four months of continuous deployment. The system currently
indexes 64,180+ chunks across a shared canonical corpus, maintains 99.97\% vector
coverage, and operates at p95 latency below one second. Three primary contributions
distinguish this work:

\begin{enumerate}

\item \textbf{Pain-weighted salience} (\S3.5). We introduce a retrieval scoring signal
that explicitly models incident severity:
$\texttt{salience} = \texttt{recency} \times \texttt{pain} \times \texttt{importance}$,
where $\texttt{pain} \in [0.1, 1.0]$ encodes the operational cost of the recorded
experience---from trivial notes (0.1) to production outages (1.0). To our knowledge, this
is the first documented retrieval signal in the memory systems literature that treats
incident severity as a first-class dimension. Preliminary ablation studies on post-incident
queries suggest this dimension contributes positively to retrieval quality (see \S5.6).

\item \textbf{Enforced shadow discipline} (\S3.6). We describe a methodology for
validating ranking changes against a shadow telemetry baseline before deployment,
implemented as an architectural constraint rather than a documented suggestion. The
mechanism is codified via environment variable gating
(\texttt{NOX\_SALIENCE\_MODE=shadow|active}), enforced by a cron-checked health endpoint,
with a mandatory minimum observation period of seven days. During the salience validation
run (Fase~1.7b-b), this produced 191 promotion candidates, 16,608 review cases, and
45,743 archive recommendations before any activation. This is, to our knowledge, the first
memory system paper to codify such a discipline as an architectural constraint rather than
a suggested practice.

\item \textbf{Shared-canonical multi-agent design} (\S3.7). Six agents share a single
\texttt{chunks} table distinguished by \texttt{source\_file} prefix rather than separate
stores. Cross-agent retrieval requires no synchronization because partitioning never
occurred. When Forge records a decision about infrastructure, Atlas retrieves it directly
on the next relevant query---not because data was replicated, but because the design
assumed shared context from the start.

\end{enumerate}

Secondary contributions include: a reproducible evaluation harness with nDCG@10, MRR,
Recall@10, and Precision@5 over 50 curated golden queries; four months of operational
evidence across five infrastructure upgrades without data loss; and comparative analysis
against strong retrieval baselines (see \S5).

\subsection{Paper Organization}

Section~2 surveys related work across knowledge graph--augmented retrieval, agent memory
systems, and production-oriented frameworks, identifying the specific gaps this work
addresses. Section~3 describes the system architecture across storage, retrieval, knowledge
graph, and operational discipline layers. Section~4 details the evaluation methodology,
including the golden query harness, baseline selection, and three-corpus design. Section~5
presents experimental results, including hybrid pipeline ablations, baseline comparisons,
pain dimension validation, and the shadow discipline case study. Section~6 discusses
deployment scope, limitations, and future work. Section~7 concludes.

% ---------------------------------------------------------------
% SECTION 2 — RELATED WORK
% ---------------------------------------------------------------
\section{Related Work}

Agent memory research has accelerated rapidly since the introduction of
retrieval-augmented generation \cite{lewis2020rag}, with recent work spanning knowledge
graph--augmented retrieval, autonomous memory management, and production deployment
frameworks. We survey the most directly relevant systems and identify the specific
dimensions where this work diverges.

\subsection{Knowledge Graph--Augmented Retrieval}

\textbf{GraphRAG} \cite{edge2024graphrag} introduced a pipeline that uses LLM extraction
to build entity-relation graphs over document corpora, applies community detection (Leiden
algorithm) to identify thematic clusters, and generates multi-level summaries for
query-focused retrieval. The approach substantially improves global query coverage---queries
requiring synthesis across many documents---compared to naive RAG. nox-mem adopts the
LLM-extraction paradigm for knowledge graph construction but diverges in objective:
GraphRAG optimizes for global summarization via community traversal, while nox-mem targets
entity-centric local retrieval for operational contexts, with closed-enum edge typing
designed to support downstream blast-radius analysis. GraphRAG does not model retrieval as
an operational concern and has no analogue to the discipline mechanisms described in
\S3.5--3.6.

\textbf{HiRAG} \cite{huang2025hirag} extends the KG-augmented paradigm with hierarchical
knowledge representation and iterative reasoning chains, demonstrating state-of-the-art
performance on multiple QA benchmarks. The hierarchical approach---traversing from coarse to
fine-grained graph levels---optimizes for accuracy on complex reasoning tasks. nox-mem
instead maintains a flat three-layer fusion pipeline (\S3.4), accepting the
accuracy-latency trade-off explicitly: p95 below one second at 64K chunks serves
operational use cases where query latency is a usability constraint, not just a benchmark
metric. HiRAG does not address multi-agent deployment, evaluation harnesses for operational
corpora, or retrieval discipline.

\textbf{HippoRAG} \cite{guo2024hipporag} draws on neurobiological memory models to improve
associative retrieval via hippocampal-inspired graph structures. The biological analogy
provides useful intuition for the salience formula introduced in \S3.5, where pain-weighted
recency mirrors the well-documented human tendency to retain high-cost experiences more
strongly than low-cost ones. HippoRAG does not formalize this dimension as a retrieval
signal.

\subsection{Agent Memory Systems}

\textbf{MemGPT} \cite{packer2023memgpt} proposed treating the LLM itself as an operating
system, with the agent autonomously managing its own context through function calls that
page between main context and external storage. The architecture gives agents direct
control over their memory---a compelling design for single-agent autonomy. nox-mem adopts
the opposite philosophy: memory is managed by an external pipeline (file watcher, ingest
router, nightly KG extraction), and agents are consumers of a shared canonical store rather
than managers of isolated state. This separation enables cross-agent retrieval as a
structural property while imposing the discipline that agents cannot corrupt their own
context.

\textbf{Mem0} \cite{chhikara2025mem0} introduced a production-focused memory layer with
LLM-driven memory extraction, deduplication, and temporal updates. Evaluated on the LOCOMO
benchmark, Mem0 demonstrates +26\% accuracy improvement over baselines on long-context
conversation scenarios. nox-mem shares the production emphasis but diverges in two
important respects. First, Mem0 relies on autonomous LLM editing decisions to update and
merge memories; nox-mem requires shadow-period validation before any retrieval-affecting
change propagates to production, specifically to mitigate silent regression risk. Second,
Mem0's multi-user design partitions memory by \texttt{user\_id}, treating agents as
user-addressable endpoints; nox-mem's shared canonical design treats agents as roles within
a single trusted operational context.

\textbf{A-MEM} \cite{xu2025amem} introduced agentic memory with Zettelkasten-inspired
structure: notes are automatically tagged, interlinked, and contextualized by an LLM
controller, enabling dynamic memory evolution. The approach handles unstructured
conversation well but does not address operational corpora with explicit schema
requirements, retention policies, or incident-severity differentiation. A-MEM has no
evaluation harness independent of LLM-generated quality judgments, and does not address
multi-agent deployments or ranking change discipline.

\subsection{Production-Oriented Frameworks}

\textbf{Cognee} \cite{topoteretes2024cognee} provides an open-source AI memory framework
implementing an Extract-Cognify-Load (ECL) pipeline with native knowledge graph support
and hybrid retrieval. Of the systems surveyed, Cognee is architecturally closest to
nox-mem: both maintain native KGs, both support hybrid retrieval, and both target
production deployment scenarios. Cognee achieves 3/5 on our architectural parity dimensions
(Table~\ref{tab:comparison}). The key differentiators are shadow discipline---Cognee provides
no enforced validation gate for retrieval-affecting changes---and the evaluation harness:
Cognee lacks a reproducible evaluation methodology with standard IR metrics over a fixed
golden set. nox-mem's eval harness (\S4) enables the kind of regression detection that the
shadow discipline gate is designed to act on.

\textbf{LangChain Memory} provides key-value and buffer memory primitives widely used in
production applications. The design prioritizes API simplicity and framework compatibility
over retrieval quality; it does not support hybrid retrieval, knowledge graph construction,
or evaluation. It serves as a practical lower bound in comparative analysis.

\subsection{Foundational Retrieval Methods}

\textbf{Retrieval-Augmented Generation} \cite{lewis2020rag} established the foundational
architecture combining parametric LLM knowledge with non-parametric document retrieval,
enabling factual question answering over dynamic corpora. nox-mem extends this foundation
to persistent, evolving operational corpora with explicit retention, pain weighting, and
multi-agent routing.

\textbf{Reciprocal Rank Fusion} \cite{cormack2009rrf} provides the rank-combination
mechanism at the core of nox-mem's hybrid retrieval layer. Cormack et al.\ demonstrated
that RRF consistently outperforms individual rank learning methods and Condorcet-based
fusion without requiring system-specific tuning. We adopt RRF with $k=60$ as the fusion
parameter, consistent with the paper's recommendations, as the final combination step over
BM25 and dense retrieval ranked lists.

\subsection{Positioning: Where nox-mem Fits}

Table~\ref{tab:comparison} summarizes the architectural dimensions across surveyed systems.
No existing system in the literature addresses all five dimensions: native knowledge graph
support, hybrid retrieval, a reproducible evaluation harness, multi-agent design, and
enforced shadow discipline. The two dimensions with zero coverage in the literature---pain-weighted
salience as a retrieval signal and shadow discipline as an architectural constraint---are
the primary novelty claims of this paper.

\begin{table}[t]
  \caption{Architectural comparison across memory systems.}
  \label{tab:comparison}
  \centering
  \small
  \begin{tabular}{lllllll}
    \toprule
    System & KG Native & Hybrid Retr. & Eval Harness & Multi-Agent & Shadow Disc. & Score \\
    \midrule
    \textbf{nox-mem (ours)} & closed-enum, 7 & FTS5+Gemini+RRF & nDCG/MRR/Recall & shared canonical & enforced $\geq$7d & \textbf{5/5} \\
    GraphRAG \cite{edge2024graphrag} & community det. & via KG queries & none & none & none & 1.5/5 \\
    MemGPT \cite{packer2023memgpt} & none & embedding-first & none & per-agent state & none & 1.5/5 \\
    Mem0 \cite{chhikara2025mem0} & optional (v2) & vector-only & LOCOMO only & user\_id partition & none & 1.5/5 \\
    A-MEM \cite{xu2025amem} & Zettelkasten & semantic-first & none & none & none & 1.0/5 \\
    HiRAG \cite{huang2025hirag} & hierarchical & multi-level & task-specific & none & none & 2.5/5 \\
    Cognee \cite{topoteretes2024cognee} & ECL pipeline & hybrid & ad-hoc & optional & none & 3.0/5 \\
    LangChain Memory & none & key-value & none & session\_id & none & 0.5/5 \\
    \bottomrule
  \end{tabular}
\end{table}

The closest system, Cognee, covers 60\% of these dimensions. The average score across the
seven surveyed systems is 1.6/5. Section~3 describes how nox-mem addresses all five, and
Sections~4--5 provide the empirical grounding for the claims in this table.

% ---------------------------------------------------------------
% SECTION 3 — SYSTEM ARCHITECTURE
% ---------------------------------------------------------------
\section{System Architecture}

nox-mem is a production memory system built on SQLite with three primary subsystems---storage,
retrieval, and knowledge graph---plus an operational discipline layer that governs how the
system evolves. The implementation is in TypeScript using \texttt{better-sqlite3} for
synchronous database access, \texttt{sqlite-vec} for approximate nearest-neighbor search,
and Gemini for embedding generation and LLM-based KG extraction. The system has operated
continuously since March 2026 with 12 schema migrations and no irrecoverable data loss (one
metadata-loss event on 2026-04-25 was recovered via re-ingestion; see \S6.2).

\subsection{System Overview}

\begin{figure}[t]
  \centering
  \includegraphics[width=\linewidth]{figures/figure1-system-overview.pdf}
  \caption{System overview. Six personas read and write through a unified interface layer
    (CLI / MCP / HTTP) into a single canonical SQLite corpus (no per-agent silos).
    Retrieval runs as a three-layer hybrid (BM25 $\to$ 3072-d Gemini cosine $\to$ RRF,
    $k=60$). The shadow-mode pipeline and append-only \texttt{ops\_audit} are cross-cutting
    governance: every ranking change ships shadow-first; every destructive op is wrapped by
    \texttt{withOpAudit()} snapshot.}
  \label{fig:overview}
\end{figure}

The entry point for all CLI operations is \texttt{dist/index.js} (the compiled package
binary). An HTTP API on port 18802 exposes search, knowledge graph, health, and operational
endpoints for agent integration.

\begin{verbatim}
Source files (.md / .docx / .pptx / .pdf / entity files)
    |
    v
ingest-router (inotifywait watcher + manual CLI)
    |
    +------> chunks table (canonical, 64K+ rows)
                |
                +------> chunks_fts (FTS5 BM25 index)
                +------> vec_chunks (sqlite-vec, 3072d embeddings)
                +------> kg_entities + kg_relations (LLM-extracted, nightly)
                +------> search_telemetry (shadow-mode + eval harness)
                |
    query --> search() pipeline:
                Layer 1: FTS5 BM25 candidates
                Layer 2: Gemini semantic candidates
                Layer 3: RRF fusion (k=60) + salience boost + section_boost
                --> ranked results (p95 < 1s @ 64K chunks)
\end{verbatim}

\subsection{Storage Layer}

The canonical store is a single SQLite database accessed via \texttt{better-sqlite3} for
synchronous, blocking reads---a deliberate choice that simplifies the concurrency model for
a single-node deployment. Three coordinated tables implement the storage layer.

\texttt{chunks} is the central table, containing content text, metadata
(\texttt{source\_file}, \texttt{chunk\_type}, \texttt{created\_at}, \texttt{last\_seen}),
and operational fields (\texttt{pain} REAL DEFAULT 0.2, \texttt{importance} REAL,
\texttt{section} TEXT, \texttt{section\_boost} REAL, \texttt{retention\_days}). The schema
has evolved through 12 versions (v1 through v12) since March 2026; all migrations have been
applied to the live database without downtime or data loss. Key schema additions include
\texttt{retention\_days} (v8, typed retention by chunk category), \texttt{pain} (v9,
severity signal), and \texttt{section} with \texttt{section\_boost} (v10, structural
weighting for entity file sections).

\texttt{chunks\_fts} is an FTS5 virtual table providing BM25-ranked full-text search over
chunk content. FTS5 vanilla AND-mode behavior---requiring all query terms to
match---means that natural language queries frequently return zero results when run against
the lexical index alone. This structural characteristic of FTS5 makes hybrid retrieval a
necessity rather than an optimization, a finding confirmed in evaluation (\S5.1).

\texttt{vec\_chunks} stores 3072-dimensional Gemini embeddings
(\texttt{gemini-embedding-001}) via the \texttt{sqlite-vec} extension, enabling
approximate cosine similarity search without external vector store infrastructure. Current
vector coverage is 99.97\% across 64,180+ chunks.

Retention policy is encoded directly in the schema via \texttt{retention\_days}: feedback
and person records are set to \texttt{NULL} (never decay); decisions and projects have
365-day retention; team context 120 days; daily notes 90 days; pending items 30 days. This
ensures that the system's own operating rules cannot be silently evicted by a generic decay
policy.

\subsection{Retrieval Layer}

The retrieval pipeline operates in three sequential layers, with results from layers one and
two merged via Reciprocal Rank Fusion \cite{cormack2009rrf}:

\textbf{Layer 1---Lexical retrieval.} FTS5 BM25 search expands the query against the
\texttt{chunks\_fts} index, returning up to $k_1$ ranked candidates. Results are
deduplicated by chunk ID before passing to fusion.

\textbf{Layer 2---Semantic retrieval.} The query is embedded using
\texttt{gemini-embedding-001} (3072d), and cosine similarity search is performed against
\texttt{vec\_chunks} via \texttt{sqlite-vec}, returning up to $k_2$ ranked candidates.

\textbf{Layer 3---RRF fusion and boosting.} Candidates from layers one and two are merged
via RRF with $k=60$. The fused score is then multiplied by two additive boost terms:
\texttt{section\_boost} (compiled section: $2.0\times$, frontmatter: $1.5\times$,
timeline: $0.8\times$, unstructured: $1.0\times$) and \texttt{salience} (\S3.5). Boosting
is implemented as additive contributions to the log-transformed RRF score to prevent
multiplicative stacking, which prior operational experience identified as a source of score
distribution collapse.

The pipeline achieves p95 search latency below one second across the 64K chunk corpus,
measured in production over the observation period. The \texttt{-{}-no-hybrid} flag disables
semantic retrieval for latency-constrained contexts where lexical precision is acceptable.

\subsection{Knowledge Graph Layer}

An LLM-extracted knowledge graph augments chunk retrieval with structured relationship
data. Entities and relations are extracted incrementally via nightly batch runs using
\texttt{gemini-2.5-flash} and stored in \texttt{kg\_entities} ($\approx$402 entities) and
\texttt{kg\_relations} ($\approx$544 relations).

Entity types follow a closed taxonomy of 10 categories: person, project, agent, tool,
concept, organization, technology, document, decision, and metric. Relation edges use
free-form \texttt{relation\_type} text combined with a closed-enum
\texttt{relation\_reason} field with seven canonical values: \texttt{depends\_on},
\texttt{derived\_from}, \texttt{opposes}, \texttt{extends}, \texttt{replaces},
\texttt{mentions}, and \texttt{unknown}. The closed enum enables downstream tooling---\texttt{impact~\textless entity\textgreater} computes one-hop blast radius with
reason-priority weighting; \texttt{detect-changes} maps git diffs to affected KG
entities---without requiring LLM inference at query time.

Edge type accuracy required explicit engineering. Initial extraction with an optional enum
field and a ``use unknown if unsure'' prompt instruction produced 86\% \texttt{unknown}
labels across n=100 sampled relations. A three-path fix---a revised prompt with explicit
examples for each non-unknown category; a code-side defensive normalization map
(\texttt{RELATION\_TYPE\_TO\_REASON}, 24 input aliases covering PT-BR and EN free-text
variants that collapse to the 7 canonical reasons); and a post-extraction validation
pass---improved the classification rate from 14\% to 56\% (4$\times$ improvement),
equivalently reducing the \texttt{unknown} rate from 86\% to 44\% (n=100 sample). This
experience reinforces a general principle: LLM-extracted structured fields with optional
enums and uncertainty-deferring prompts will systematically underperform; defensive
code-side normalization is required.

\subsection{Pain-Weighted Salience}

The salience formula is:

\begin{equation}
  \texttt{salience}(\mathit{chunk}) =
    \texttt{recency}(\mathit{chunk})
    \times \texttt{pain}(\mathit{chunk})
    \times \texttt{importance}(\mathit{chunk})
  \label{eq:salience}
\end{equation}

\noindent where:
\begin{itemize}
  \item $\texttt{recency} \in [0,\,1]$ is an exponential decay function over the
    \texttt{last\_seen} timestamp, parameterized by the chunk's \texttt{retention\_days}
    value;
  \item $\texttt{pain} \in [0.1,\,1.0]$ encodes the operational severity of the recorded
    experience, annotated via \texttt{<!-- pain: 0.X -->} markers in entity files. Values
    are assigned on a five-point semantic scale: 0.1 (trivial note), 0.3 (minor friction),
    0.5 (significant issue), 0.7 (major incident), 1.0 (production outage with data loss
    or downtime);
  \item $\texttt{importance} \in [0,\,1]$ is derived from \texttt{mention\_count} across
    the corpus and entity-type priors.
\end{itemize}

The formula is multiplicative: a high-pain chunk with low recency can still score above a
low-pain chunk with high recency when the severity differential is large. This mirrors the
asymmetry in human operational memory, where costly experiences are retained and surfaced
preferentially over routine recent ones \cite{guo2024hipporag}.

The salience signal is exposed in shadow mode at \texttt{/api/health.salience} without
affecting production ranking, enabling comparison against the baseline before activation
(\S3.6).

\subsection{Enforced Shadow Discipline}

\begin{figure}[t]
  \centering
  \includegraphics[width=0.9\linewidth]{figures/figure3-shadow-state-machine.pdf}
  \caption{Shadow discipline state machine. Every ranking-affecting change traverses this
    machine. Activation requires three independent signals: at least seven days of shadow
    telemetry, a healthy distribution (no degenerate clustering), and explicit human
    approval. Regression at any point reverts via a single environment-variable
    flip---zero schema rollback required.}
  \label{fig:shadow}
\end{figure}

The shadow discipline mechanism enforces a separation between observing a retrieval change
and deploying it. Any modification that affects ranking---scoring formula changes, boost
adjustments, new retrieval layers---must run in shadow mode for a minimum of seven days
before activation. This is an architectural constraint, not a documented recommendation.

Implementation relies on three components: (1)~an environment variable gate
(\texttt{NOX\_SALIENCE\_MODE=shadow|active}) that controls whether the salience score
affects production ranking or is only written to \texttt{search\_telemetry}; (2)~a health
endpoint (\texttt{/api/health.opsAudit}) that cron checks every 15 minutes and alerts via
Discord if the shadow observation period has not been met; (3)~an append-only audit log
(\texttt{ops\_audit} table) with database triggers that block DELETE and UPDATE on rows
with terminal status, ensuring the shadow-period record cannot be retroactively modified.

The distinction from A/B testing is important: A/B testing requires live traffic over a
production population. Shadow discipline for operational/personal corpora has no such
traffic---queries are infrequent, irregular, and non-i.i.d. The shadow mechanism instead
collects telemetry on what the proposed ranking \emph{would have done} on each live query,
accumulating promotion, review, and archive candidate distributions across the observation
period. The distribution comparison then drives the activation decision.

The salience validation run (Fase~1.7b-b) ran for seven days before activation, producing
191 promotion candidates, 16,608 review cases, and 45,743 archive recommendations. The
distribution was reviewed, found consistent with design intent, and activation proceeded.
Counterfactual analysis suggests that the April~25 reindex incident---the motivating example
in \S1---would have been detected during a shadow period: the reindex operation altered
chunk distributions in ways that would have surfaced as anomalous in telemetry within hours.

\subsection{Shared-Canonical Multi-Agent Design}

Six agents (Atlas, Boris, Cipher, Forge, Lex, and Nox) share a single \texttt{chunks}
table. Agent-specific content is distinguished by \texttt{source\_file} path prefix
(\texttt{agents/<name>/...}) rather than separate tables, schemas, or databases.
Cross-agent queries use \texttt{cross-search} and \texttt{cross-kg} CLI commands, which
apply additional \texttt{source\_file} filters to retrieve context generated by a specific
agent.

The design assumes a single trusted operator. It is explicitly not suitable for
multi-tenant SaaS deployments, where user isolation is a security requirement rather than
an engineering choice. Within the trusted-operator constraint, shared-canonical design
provides cross-agent intelligence as a structural property: when Forge ingests a decision
about an infrastructure change, Atlas retrieves it on the next architecturally relevant
query without any synchronization step.

The operational tooling layer---\texttt{detect-changes}, \texttt{impact},
\texttt{api-impact}, \texttt{consolidate-merge}, \texttt{kg-reclassify}---extends the
system to support multi-agent workflows beyond retrieval, enabling agents to query the KG
for blast-radius analysis before executing changes, and to surface merge candidates across
agent-sourced entities.

\fi % end \iffalse — monolith §1-3 suppressed

% ---------------------------------------------------------------
% SECTIONS 4–7 + APPENDICES — sourced from sec_4_7.tex
% ---------------------------------------------------------------
% sec_4_7.tex — Sections 4–7 + Appendices A–D for arXiv submission
% Source: paper-draft-sec4-7.md (converted 2026-05-04)
% Includes confirmed experimental data: R01b, BM25 Pyserini, E10 pain ablation,
%   multilingual-e5-base, LOCOMO FTS5 baseline, cross-agent storage quantification.
% Language: English (paper language)

% ---------------------------------------------------------------
% SECTION 4 — METHODS
% ---------------------------------------------------------------
\section{Methods}

We describe the evaluation framework, shadow-mode methodology, and calibration procedures
used to validate the three primary contributions. All experiments were pre-registered in
\texttt{03-experiments-needed.md} before results were collected, following the
open-evaluation norm advocated by \cite{rogers2021just}.

\subsection{Evaluation Harness}

Our primary evaluation uses \textbf{nDCG@10}, \textbf{MRR} (Mean Reciprocal Rank),
\textbf{Recall@10}, and \textbf{Precision@5} computed over a set of 60 internally curated
golden queries (dataset R01c-v1.1, post-cure refresh as of 2026-05-06; see \S5.1 ``Gold standard revision note''). These metrics follow the
standard IR evaluation methodology described in \cite{manning2008introduction}.

\textbf{Golden query construction.} Queries span eight categories reflecting the
operational nature of the corpus: \texttt{entity} (specific named entities---agents, tools,
decisions), \texttt{procedure} (how-to operational steps), \texttt{concept} (abstract
architectural notions), \texttt{security} (vulnerability and mitigation queries),
\texttt{decision} (architectural choices and their rationale), \texttt{cross-agent}
(questions whose answer originates from a different agent's memory space),
\texttt{temporal} (time-anchored recall, e.g., ``what changed in late April''), and
\texttt{negative} (6 queries, 12\% of set, for which the correct answer is that no
relevant chunk exists---testing specificity against hallucination risk). Each query was
authored by the single curator with a relevance label set ($0 =$ not relevant,
$1 =$ partially relevant, $2 =$ highly relevant) over the top-20 retrieved candidates.

\textbf{Held-out subset (R01c).} An additional ten queries (R01c, Q51--Q60) are designated
held-out: they were locked before any retrieval tuning and are evaluated only once per
major system revision, functioning as a proxy for external-curator independence. The
total golden-query budget is therefore 60 (50 main R01b $\cup$ 10 held-out R01c).
Performance on R01c is reported separately from the 50-query main set R01b to avoid
optimistic bias from iterative query refinement.

\textbf{Internal-curator bias mitigation.} We acknowledge that golden queries authored by
the same individual who built the system introduce construct validity risk. Three
mitigations are applied: (i)~the held-out R01c subset was frozen before the final tuning
sprint; (ii)~external corpora (BEIR TREC-COVID, Stack Exchange---\S5.2) use third-party
curated relevance judgments; and (iii)~six negative queries test the boundary condition
most susceptible to self-serving bias. As a fourth mitigation, we evaluate nox-mem against
10 queries authored by NIST professional assessors (TREC-COVID Round 5
\cite{thakur2021beir}), selected via TF-IDF $k$-means clustering ($k$=10, seed=42) over
the 50 canonical BEIR topics to maximise lexical diversity (avg pairwise Jaccard = 0.097).
Vocabulary overlap with the internal golden-50 set was 0.0\%, confirming the two sets
probe complementary terminology. See \S5.2 for cross-corpus generalization results.

\subsection{Shadow-Mode Methodology}

Any change that affects retrieval ranking in the production system is subject to mandatory
shadow validation before activation. The protocol is enforced architecturally via the
environment variable \texttt{NOX\_SALIENCE\_MODE}, which accepts three values:
\texttt{shadow} (collect both old and new scores in \texttt{search\_telemetry} without
applying the new ranking), \texttt{active} (apply new ranking), and \texttt{off} (disable
the feature entirely).

\textbf{Telemetry collection.} In shadow mode, every search call writes a row to
\texttt{search\_telemetry} containing: \texttt{query\_text} (opt-in,
\texttt{NOX\_SEARCH\_LOG\_TEXT=1}), \texttt{old\_score}, \texttt{new\_score},
\texttt{top\_chunk\_ids}, and \texttt{top\_scores}. This enables offline comparison of
old and new score distributions without exposing users to the changed ranking.

\textbf{Activation gate.} Shadow validation runs for a minimum of seven calendar days.
After the shadow period, the stored distribution is analyzed: if the new score distribution
shows statistically meaningful separation from the old distribution (inspected via
percentile comparison and visual histogram), and if no ranking inversion is detected on a
manually reviewed 10-query spot check, the feature advances to \texttt{active}. The
seven-day minimum is not a guideline---it is a hard constraint codified in the cron
configuration that governs feature activation. This design choice is motivated by the
incident of 2026-04-25, where a ranking-affecting change reached production without any
offline validation period, causing 183 entity records to lose their structured metadata
without triggering any alert (\S6.2).

\textbf{Case study: Fase~1.7b-b salience activation.} During the seven-day shadow period
for the pain-weighted salience formula, the system collected telemetry over 191 promotion
candidates, 16,608 review candidates, and 45,743 archive candidates. The distribution
separated clearly across all three tiers. Only after this distribution analysis did we
advance \texttt{NOX\_SALIENCE\_MODE} from \texttt{shadow} to \texttt{active}. This case
study is documented in detail in Appendix~B.

\subsection{Pain Weighting Calibration}

The \texttt{pain} field is a real-valued annotation in $[0.1,\,1.0]$ attached to each
chunk at ingest time. Annotation is currently manual, using the \texttt{pain: X.X} marker
syntax in entity files, and defaults to 0.2 for unannotated content.

\textbf{Calibration heuristics.} Based on approximately three months of operational experience, the
following calibration anchors were established: 0.1 (trivial notes, meeting summaries with
no operational consequence); 0.2 (default, documentation and informational content);
0.3--0.4 (decisions with moderate reversibility risk); 0.5--0.7 (production incidents with
bounded impact---recoverable within one session); 0.8--0.9 (incidents causing data loss or
multi-hour outages); 1.0 (catastrophic incidents---unrecoverable data loss, multi-day
downtime, or security breach). The calibration is conservative: in ambiguous cases,
annotators are instructed to use the lower bound of the relevant range, then escalate only
if post-incident analysis reveals higher severity.

\textbf{Engineering rationale for scale and aggregation form.} The five-point scale and
the $10\times$ spread between extreme values are engineering choices, not psychometric or
biologically grounded measurements. The scale structure follows established incident
management taxonomies \cite{pagerduty2023severity,beyer2016site}: production operations
teams routinely discriminate between severity tiers (e.g., P1 outage vs.\ P5
informational) and dispatch resources accordingly. Multiplicative aggregation over
$\texttt{recency} \times \texttt{pain} \times \texttt{importance}$ is preferred over
additive because an additive offset shrinks relative to RRF scores as corpus size grows,
whereas multiplicative coupling preserves the severity ratio in log-scale ranking
(see \S3.3). We explicitly acknowledge that these choices have not been empirically
ablated; the paper fixes this calibration and measures retrieval performance under it.
Ablation across spread values and aggregation forms is documented as future work in \S6.3.

\textbf{Pain dimension used in the salience formula.} See Equation~\ref{eq:salience},
where $\texttt{recency} \in [0,1]$ is an exponential decay over \texttt{last\_seen}
timestamp, $\texttt{pain} \in [0.1,1.0]$ is the manual annotation, and
$\texttt{importance} \in [0,1]$ is derived from \texttt{mention\_count} and
\texttt{entity\_type} prior. The multiplicative structure means that a high-pain chunk
remains salient even as its recency decays---which is the core behavioral claim of
Contribution~1.

\textbf{Annotation coverage.} The corpus contains 61,257 chunks; pain annotation is
currently applied selectively to chunks derived from incident entity files. Future work
includes LLM-driven automatic pain classification over the full corpus (\S6.5).

\subsection{Edge Typing Extraction}

\begin{figure}[!ht]
  \centering
  \includegraphics[width=\textwidth, keepaspectratio]{figures/figure4-kg-edge-typing.pdf}
  \caption{KG edge typing flow. Gemini emits free-text relation labels; a defensive
    three-path normalizer collapses them into a closed enum of seven canonical reasons
    enforced by a schema-v12 \texttt{CHECK} constraint. Before fix: 14\% classified,
    86\% unknown. After fix: 56\% classified, $4\times$ improvement ($n$=100).}
  \label{fig:edgetyping}
\end{figure}
\FloatBarrier

The knowledge graph relation schema uses a closed-enum field \texttt{relation\_reason}
with seven values: \texttt{depends\_on}, \texttt{derived\_from}, \texttt{opposes},
\texttt{extends}, \texttt{replaces}, \texttt{mentions}, and \texttt{unknown}. The goal of
edge typing is to enable blast-radius queries (e.g., ``what does component X depend on?'')
that are impossible with untyped relations.

\textbf{Prompt design and the unknown-rate problem.} An initial prompt that marked
\texttt{relation\_reason} as an optional field with the instruction ``use \texttt{unknown}
if unsure'' produced 86\% unknown-typed relations across $n$=100 sampled extractions,
rendering the typed KG practically useless for blast-radius queries. The fix applied a
three-path defensive normalization strategy: (i)~a revised prompt that provides explicit
examples for each of the six non-unknown categories and makes \texttt{unknown} a last
resort rather than a default; (ii)~a code-side defensive map
(\texttt{RELATION\_TYPE\_TO\_REASON}, 24 entries) that normalizes LLM-produced free-text
variants---including PT-BR and EN aliases---to the 7 canonical enum values; and (iii)~a
post-extraction validation pass that re-prompts any row where the LLM output did not match
the closed enum. After this fix, the enum coverage rate improved from 14\% to 56\% on
$n$=100 sampled relations ($4\times$ improvement); equivalently, the \texttt{unknown} rate
decreased from 86\% to 44\%. This is a self-reported coverage rate: no human-annotated
ground truth was used and no inter-rater agreement study (Cohen's $\kappa$) was conducted
(see \S6.3 for the full limitation statement).

\textbf{Current KG state (2026-05-03).} The production graph contains approximately 402
entities and 544 relations, extracted incrementally by a nightly Gemini 2.5 Flash job. KG
extraction uses the full \texttt{gemini-2.5-flash} model (not the lite variant) given the
low daily volume and the higher extraction quality requirements.

\FloatBarrier
\subsection{Statistical Methodology}

All retrieval experiments report \textbf{3-run mean $\pm$ standard deviation} with Bessel
correction. The system is operationally deterministic for identical queries against a
static corpus (no stochastic ranking), so run-to-run variance arises primarily from corpus
index state and warm-cache effects; empirically, standard deviation across runs is
consistently below 0.001 for nDCG@10.

For small-N validations---specifically the pain dimension experiment (E10, $n=31$
post-incident queries)---we additionally report \textbf{bootstrap 95\% confidence
intervals} computed with 10,000 resamples. Given the small sample, bootstrap CI is the
appropriate uncertainty quantification; we do not use asymptotic normal approximations for
$n < 30$.

Effect sizes for ablation experiments (\S5.3) are reported as absolute
$\Delta$~nDCG@10 rather than relative percentages, following the recommendation of
\cite{fuhr2018some} to avoid inflating small absolute differences through percentage
framing.

The transition to \S5 follows directly: the methods described above define the evaluation
apparatus; the next section applies that apparatus to produce results across five
experimental questions.

% ---------------------------------------------------------------
% SECTION 5 — EXPERIMENTS AND RESULTS
% ---------------------------------------------------------------
\needspace{6\baselineskip}
\section{Experiments and Results}

We report results across five experimental themes: (5.1) internal corpus baselines
including external lexical/dense comparisons (R01a/b/c, E1--E3) establishing the
hybrid pipeline as a minimum viable requirement; (5.2) cross-corpus generalization
on BEIR TREC-COVID and LOCOMO (E4+E5+E11); (5.3) ablation studies isolating each
architectural layer (E6--E9); and (5.4--5.6) targeted validation of the two novel
contributions---pain weighting (E10) and cross-agent intelligence (E12). All pre-registered
hypotheses from \texttt{03-experiments-needed.md} are stated before results; experiments
flagged for a follow-up revision are noted explicitly in their respective tables (e.g.,
the three layer ablations E7--E9 in Table~\ref{tab:ablation}). Golden queries
pre-registered at \S3.8 ensure post-hoc construction bias is bounded.

\subsection{Internal Corpus Baselines and External Comparisons (R01a/b/c, E1--E3)}

\textbf{Pre-registered hypothesis (R01a):} The hybrid pipeline will outperform FTS-only
BM25 by a substantial margin on natural-language queries over the operational corpus.

\textbf{Result (confirmed, R01b/R01c, 2026-05-03; gold standard refresh R01c-v1.1, 2026-05-06):} Table~\ref{tab:hybrid} shows the
primary comparison. FTS5 vanilla BM25 achieves nDCG@10 $= 0.0000$ (exact zero on the post-cure $n$=60 corpus) on
natural-language queries against the operational corpus. This is not an artifact of query
phrasing: FTS5 applies AND-strict matching by default, which means any multi-word
natural-language query that does not appear verbatim in the corpus returns zero results.
This is a structural property of the retrieval system, not a tuning failure, and it
confirms that hybrid retrieval is a minimum viable requirement rather than an optimization
for this corpus type. nox-mem hybrid achieves a $4.0\times$ lift over the strongest
external lexical baseline (BM25 Pyserini, Anserini-tuned).

\textbf{Gold standard revision note (v1.0 $\to$ v1.1).} Between paper draft v1.0 (R01b, $n$=50) and v1.1 (R01c-v1.1, $n$=60), the gold standard was revised in three ways: (i) 11 queries with empty \texttt{expected\_chunk\_ids} (questions for which no chunk in the corpus actually answers; documentation gaps) were moved to \texttt{category=negative} so they correctly score 0.000 in any system, instead of distorting non-negative category means with false zeros; (ii) 3 orphan IDs that referenced chunks deleted by re-ingest events were updated to current equivalents; (iii) 3 temporal queries (Q70/Q87/Q88) were extended with newly available timeline events. The revision strengthens the absolute $\Delta$ (0.509 $\to$ 0.583 nDCG@10) without changing the qualitative finding. Both v1.0 and v1.1 numbers are kept in the repository for transparency. v1.1 is the canonical baseline going forward.

\begin{table}[!ht]
  \caption{Internal corpus comparison (Corpus~A, golden queries). nox-mem hybrid 3-run mean
    is on the post-cure $n$=60 set (R01c-v1.1, 2026-05-06; \S5.1 ``Gold standard revision note''); external baselines (BM25 Pyserini, multilingual-e5-base)
    are one-shot evaluations on the same $n$=60 set. FTS5 vanilla and
    nox-mem hybrid columns are directly comparable ($n$=60, R01c-v1.1, 3-run mean $\pm$ std).}
  \label{tab:hybrid}
  \centering
  \small
  \resizebox{\textwidth}{!}{%
  \begin{tabular}{lrrrrr}
    \toprule
    System & $n$ & nDCG@10 & MRR & Recall@10 & P@5 \\
    \midrule
    FTS5 vanilla (SQLite BM25, this work) & 60 & $0.0000 \pm 0.0000$ & $0.0000 \pm 0.0000$ & $0.0000 \pm 0.0000$ & $0.0000 \pm 0.0000$ \\
    BM25 Pyserini (Anserini-tuned $k_1{=}0.9$, $b{=}0.4$) \cite{yang2018anserini} & 60 & $0.1475$ & $0.1549$ & $0.2083$ & $0.0600$ \\
    multilingual-e5-base ($768$d) \cite{wang2023improving} & 60 & $0.3070$ & $0.3720$ & $0.3708$ & $0.1067$ \\
    \textbf{nox-mem hybrid (FTS5 + Gemini 3072d + RRF, $k$=60)} & 60 & $\mathbf{0.5831 \pm 0.0046}$ & $0.5445 \pm 0.0068$ & $0.7667 \pm 0.0000$ & $0.2678 \pm 0.0038$ \\
    \midrule
    $\Delta$ (hybrid $-$ BM25 Pyserini) & --- & \textbf{$+0.4356$ ($4.0\times$)} & --- & --- & --- \\
    $\Delta$ (hybrid $-$ FTS5 vanilla)  & --- & \textbf{$+0.5831$ ($+58.3$~pp)} & --- & --- & --- \\
    \bottomrule
  \end{tabular}}
  \begin{minipage}{\linewidth}
    \vspace{4pt}\footnotesize
    nox-mem hybrid 3-run mean $\pm$ std (Runs \#30/\#31/\#32, R01c-v1.1, 2026-05-07; E13 temporal-boost and E05b reason-boost held \texttt{off} to measure the core hybrid pipeline). External baselines from
    \texttt{results/E01-BM25-baseline-summary.md} (BM25 Pyserini) and
    \texttt{results/E02-E5-multilingual-baseline-summary.md} (multilingual-e5-base), one-shot on the same $n$=60 corpus.
    For comparability with paper draft v1.0, the pre-revision numbers ($n$=50, hybrid $0.5213\pm0.0004$, FTS5 $0.0123\pm0.0000$, $\Delta=+0.5090$) remain available in
    the v1.0.0 git tag (\texttt{git checkout v1.0.0} reproduces the full v1.0 paper).
    The $4.0\times$ lift over BM25 Pyserini is the cited \emph{strongest external lexical
    baseline} comparison; cross-corpus generalization (BEIR TREC-COVID, LOCOMO) is
    addressed in \S5.2 (Table~\ref{tab:beir}, Table~\ref{tab:locomo}).
  \end{minipage}
\end{table}

\subsection{Cross-Corpus Generalization (E4+E5+E11)}
\label{sec:cross-corpus}

\textbf{Pre-registered hypothesis (E4+E5+E11):} Lexical retrieval difficulty is corpus-dependent;
FTS5/BM25 performance varies by orders of magnitude across corpora with different vocabulary
overlap and topic structure. Two third-party benchmarks are evaluated: BEIR TREC-COVID
(NIST-curated relevance judgments, 171K documents) and LOCOMO (Maharana et al.\ 2024,
conversational memory). Both use frozen relevance judgments authored independently of
this work, addressing the single-curator construct-validity threat noted in \S4.1.

\begin{table}[!ht]
  \caption{Cross-corpus nDCG@10 --- BEIR TREC-COVID (171K docs, 50 queries, NIST Round 5 qrels \cite{thakur2021beir}). Cross-corpus FTS5 row is a lower bound: nox-mem hybrid was not evaluated on BEIR (requires separate ingest pipeline).}
  \label{tab:beir}
  \centering
  \small
  \resizebox{\textwidth}{!}{%
  \begin{tabular}{lrrrr}
    \toprule
    System & nDCG@10 & MRR & Recall@10 & P@5 \\
    \midrule
    BM25 (FTS5, this work) & 0.1007 & 0.2100 & 0.0017 & 0.1040 \\
    multilingual-e5-base \cite{wang2023improving} & 0.8335 & 0.8950 & 0.0212 & 0.8560 \\
    \midrule
    Cross-corpus FTS5 (this work, lower bound) & 0.2810 & --- & --- & --- \\
    \bottomrule
  \end{tabular}}
  \begin{minipage}{\linewidth}
    \vspace{4pt}\footnotesize
    BM25 uses FTS5 (SQLite) as the BM25 engine, not Pyserini/Anserini; results are directly comparable for lexical retrieval quality. nox-mem hybrid pipeline not evaluated on BEIR (would require a dedicated TREC-COVID ingest pipeline). Cross-corpus FTS5 lower bound is from the LOCOMO evaluation (Table~\ref{tab:locomo}), included for orientation only.
  \end{minipage}
\end{table}

\begin{table}[t]
  \caption{Cross-corpus nDCG@10 --- LOCOMO conversational memory benchmark (Maharana
    et al.\ 2024, 5{,}882 dialogue turns from 10 long conversations, $n$=100 stratified
    subset, seed=42).}
  \label{tab:locomo}
  \centering
  \resizebox{\textwidth}{!}{%
\begin{tabular}{lllll}
    \toprule
    System & nDCG@10 & MRR & Recall@10 & P@5 \\
    \midrule
    FTS5 (SQLite, BM25 default) & 0.2810 & 0.2795 & 0.3792 & 0.0780 \\
    BM25 (Pyserini) \cite{yang2018anserini} & \multicolumn{4}{l}{\textsc{[Deferred: future work]}} \\
    multilingual-e5-base \cite{wang2023improving} & \multicolumn{4}{l}{\textsc{[Deferred: future work]}} \\
    \textbf{nox-mem hybrid} & \multicolumn{4}{l}{\textsc{[Deferred: requires LOCOMO chunks ingest]}} \\
    \bottomrule
  \end{tabular}}
  \begin{minipage}{\linewidth}
    \vspace{4pt}\footnotesize
    LOCOMO is released under CC BY-NC 4.0 by SNAP Research \cite{maharana2024locomo}.
    Per-category breakdown ($n$=20 each, seed=42): single-hop 0.118, multi-hop 0.371,
    temporal 0.289, open-domain 0.375, adversarial 0.253. The cross-corpus FTS5 ratio ---
    LOCOMO 0.281 vs.\ internal corpus 0.0000 (post-cure $n$=60; FTS5 AND-strict on natural-language queries returns zero matches by construction) --- confirms that
    lexical retrieval difficulty is corpus-dependent. Reproducible in $<$10\,s via
    \texttt{python3 paper/publication/baselines/locomo\_eval.py full} (stdlib-only).
  \end{minipage}
\end{table}

\textbf{Cross-corpus result.} On BEIR TREC-COVID, multilingual-e5-base achieves
nDCG@10 $= 0.8335$ ($n$=50), three orders of magnitude above its internal-corpus
performance ($0.3070$, see \S5.1); the FTS5 BM25 baseline reaches $0.1007$ on TREC-COVID
versus exactly $0.0000$ on the internal corpus (post-cure $n$=60). On LOCOMO, FTS5 reaches $0.2810$ ($n$=100, well above the degenerate internal-corpus baseline).
Both lexical-baseline gaps confirm that the near-zero FTS5 result on the operational
corpus (\S5.1) is a property of \emph{this corpus} (high natural-language phrasing
density, low term overlap with stored chunks), not a general FTS5 limitation. The
nox-mem hybrid pipeline was not evaluated on either external corpus---both would
require a dedicated ingest pipeline---so the headline contribution is the
\emph{architectural pattern} (lexical $\to$ dense $\to$ RRF), not a competitive
benchmark number on third-party datasets. This is an explicit limitation
(\S6.5: deferred).

\subsection{Ablation Studies (E6--E9)}

\textbf{Pre-registered hypothesis (E6--E9):} Each of the four architectural layers (FTS5
lexical, Gemini semantic embeddings, RRF fusion, section boost) contributes positively to
nDCG@10, with each layer's removal causing a $\Delta \geq 0.03$ decrease.

\begin{table}[t]
  \caption{Layer ablation on internal corpus, Corpus~A ($n$=60 golden queries, 3-run mean
    $\pm$ std, R01c-v1.1 post-cure baseline). Two layer ablations are reported here; three additional ablations
    (no-RRF score-concat, no-salience, no-section-boost) are deferred to a follow-up
    revision (\S6.5).}
  \label{tab:ablation}
  \centering
  \resizebox{\textwidth}{!}{%
\begin{tabular}{llll}
    \toprule
    Configuration & nDCG@10 & $\Delta$ vs full hybrid & MRR \\
    \midrule
    Full hybrid (baseline)                   & $\mathbf{0.5831 \pm 0.0046}$ & ---       & $0.5445 \pm 0.0068$ \\
    FTS-only (no semantic, no RRF)           & $0.0000 \pm 0.0000$          & $-0.583$  & $0.0000 \pm 0.0000$ \\
    Hybrid without RRF fusion (E7)           & $0.5541 \pm 0.0315$          & $-0.029$  & $0.5169 \pm 0.0334$ \\
    Hybrid without salience boost (E8)       & $0.5418 \pm 0.0021$          & $-0.041$  & $0.5078 \pm 0.0030$ \\
    Hybrid without section boost (E9)        & $0.5515 \pm 0.0012$          & $-0.032$  & $0.5266 \pm 0.0015$ \\
    \bottomrule
  \end{tabular}}
  \begin{minipage}{\linewidth}
    \vspace{4pt}\footnotesize
    All ablations measured 2026-05-07 on R01c-v1.1 ($n$=60), 3-run mean $\pm$ std with E13 temporal-boost and E05b reason-boost held off.
    \textbf{Result.} All three pre-registered layer ablations confirm the hypothesis (each layer contributes $\Delta \geq 0.029$ nDCG@10):
    (E7) RRF fusion contributes $0.029$~pp; (E8) pain-weighted salience contributes $0.041$~pp (largest single-layer contribution; consistent with the E10 pain ablation in \S5.5);
    (E9) section boost contributes $0.032$~pp. Together with the FTS-only ablation ($-0.583$ vs full hybrid), this confirms the four-layer architecture is non-redundant: removing any layer measurably degrades retrieval. Toggles are environment variables (\texttt{NOX\_RRF\_DISABLE=1}, \texttt{NOX\_SALIENCE\_MODE=off}, \texttt{NOX\_SECTION\_BOOST\_MODE=off}) documented in the repository for third-party reproduction. E7 std=0.0315 is higher than E8/E9; this reflects the score-concat fallback's sensitivity to within-corpus rank distribution rather than reranker-stage non-determinism.
  \end{minipage}
\end{table}

\subsection{Pain Dimension: Empirical Ablation (E10)}
\label{sec:pain-ablation}

\textbf{Pre-registered hypothesis (E10):} On a subset of post-incident golden queries
(queries where the ground-truth answer is a chunk describing a production incident or
costly operational lesson), pain-aware retrieval (current default) will outperform
pain-uniform retrieval (\texttt{pain}=1.0 for all chunks) by $\Delta$ nDCG@10 $\geq 0.05$.

The pain-uniform counterfactual collapses all chunks to the same pain weight, effectively
reducing the salience formula to
$\texttt{salience} = \texttt{recency} \times \texttt{importance}$. This tests whether the
pain dimension adds independent retrieval signal beyond recency and importance.

\subsubsection{Methodology}

Two temporary read-only database snapshots were prepared: \texttt{pain\_real} (production
pain values, $\in [0.1, 1.0]$) and \texttt{pain\_uniform} (all chunks set to
\texttt{pain}=1.0). Both snapshots reside at \texttt{/root/.openclaw/paper-experiments/}
and were not derived from nor applied to the production database. Hybrid retrieval (FTS5
BM25 + Gemini 3072-dimensional embeddings + RRF $k=60$) was evaluated identically against
both, over $n=31$ post-incident queries. Bootstrap 95\% confidence intervals were computed
with 10,000 resamples, seed=42, over the per-query $\Delta$ nDCG@10 values.

Source: \texttt{paper/publication/baselines/pain\_ablation\_hybrid.py}; results archived in
\texttt{paper/publication/results/E10-pain-ablation-hybrid-results.md}.

\subsubsection{Aggregate Results}

\begin{table}[t]
  \caption{Pain ablation --- hybrid retrieval ($n$=31 post-incident queries, 2026-05-04).}
  \label{tab:painablation}
  \centering
  {\small
  \setlength{\tabcolsep}{4pt}
\begin{tabular}{llll}
    \toprule
    Configuration & Mean nDCG@10 & $n$ & Notes \\
    \midrule
    pain\_real (production values, $\in [0.1, 1.0]$) & \textbf{0.4469} & 31 & Confirmed --- read-only snapshot \\
    pain\_uniform (all chunks pain=1.0) & \textbf{0.4404} & 31 & Confirmed --- read-only snapshot \\
    $\Delta$ (pain\_real $-$ pain\_uniform) & \textbf{+0.0065} & --- & Directional \\
    95\% CI (bootstrap, 10,000 resamples, seed=42) & \textbf{[$-$0.0143, $+$0.0338]} & --- & CI includes zero \\
    Queries improved ($\Delta > 0$) & 1 / 31 & --- & Q55 only \\
    Queries degraded ($\Delta < 0$) & 1 / 31 & --- & Q75 only \\
    Queries unchanged ($\Delta = 0$) & 29 / 31 & --- & Gemini semantic dominates \\
    Pre-registered threshold & $\Delta \geq 0.05$ & --- & \textbf{NOT MET} \\
    \bottomrule
  \end{tabular}}
\end{table}

\textbf{Verdict: DIRECTIONAL, NOT SIGNIFICANT.} $\Delta = +0.0065$ is positive but below
the pre-registered threshold of 0.05, and the 95\% CI $[-0.0143, +0.0338]$ does not
exclude zero.

\subsubsection{Interpretation}

The pain dimension shows directional but statistically non-significant aggregate effect on
$n=31$ post-incident queries under hybrid retrieval. Disaggregation reveals the mechanism:

\begin{itemize}
  \item \textbf{29/31 queries: $\Delta = 0.000$.} Gemini semantic similarity (3072-dimensional
    cosine) produces consistent top-10 orderings that are entirely pain-insensitive. When
    the semantic model assigns clearly differentiated scores to candidates, the pain
    multiplier does not alter rank order.
  \item \textbf{Q55 (atomic pre-op backup procedure): $\Delta = +0.349$.} Pain successfully
    elevates the correct chunk from rank 2 to rank 1 when two semantically similar chunks
    receive near-identical Gemini scores.
  \item \textbf{Q75 (commit secrets rule): $\Delta = -0.148$.} Under pain\_uniform,
    FTS5 tie-breaking accidentally promotes a partially relevant security chunk more than
    under the real pain distribution. This is an artifact of the FTS-lexical component,
    not a failure of the pain signal itself.
\end{itemize}

This pattern is consistent with prior work showing that dense retrievers dominate sparse
lexical signals in fused retrieval \cite{thakur2021beir}: pain only matters in the narrow
regime where semantic scores tie, which occurred in 1 of 31 queries in this evaluation.

\subsubsection{Case Study: Q55 --- High-Pain Backup Procedure}

\textbf{Query:} ``como fazer backup pre-op atomico'' (how to perform atomic pre-op backup)

\textbf{Expected gold chunks:} ids 116179 (session handoff with backup procedure), 116380
(gateway resilience plan with backup steps).

\begin{table}[h]
  \caption{Q55 pain\_real retrieval --- nDCG@10 = 1.000.}
  \label{tab:q55real}
  \centering
  \resizebox{\textwidth}{!}{%
\begin{tabular}{llllll}
    \toprule
    Rank & Chunk ID & Score & Source & Pain & Result \\
    \midrule
    1 & 116179 & 16.39 & memoria-nox/handoffs/2026-04-21 & high & GOLD \\
    2 & 116380 & 15.87 & memoria-nox/plans/2026-04-20 & high & GOLD \\
    3 & 147900 & 15.38 & specs/202x --- archive & low & non-gold \\
    \bottomrule
  \end{tabular}}
\end{table}

\begin{table}[h]
  \caption{Q55 pain\_uniform retrieval --- nDCG@10 = 0.651.}
  \label{tab:q55uniform}
  \centering
  \resizebox{\textwidth}{!}{%
\begin{tabular}{llllll}
    \toprule
    Rank & Chunk ID & Score & Source & Pain (effective) & Result \\
    \midrule
    1 & 116179 & 16.39 & handoff & 1.0 (uniform) & GOLD \\
    2 & 147900 & 15.63 & archive spec & 1.0 (uniform) & non-gold \\
    3 & 116380 & 15.38 & resilience plan & 1.0 (uniform) & GOLD \\
    \bottomrule
  \end{tabular}}
\end{table}

\textbf{Note on score equality at rank 1.} The score 16.39 for chunk 116179 is identical
under \texttt{pain\_real} and \texttt{pain\_uniform} because 116179 (an incident-response
session handoff) already carries \texttt{pain}=1.0 in the production distribution; the
multiplicative term $\texttt{pain}=1.0$ is invariant under the two configurations.
Differences appear at ranks 2 and 3, where chunks with \texttt{pain}<1.0 (e.g., chunk
147900 archive spec) have their fused scores raised under \texttt{pain\_uniform},
displacing the second incident-derived gold chunk 116380.

\textbf{Interpretation.} When all chunks receive \texttt{pain}=1.0, the archive spec's
marginally higher semantic score is sufficient to displace 116380 from rank 2---a net rank
degradation. Pain correctly weights the incident-derived chunks above generic documentation
in the tied-score regime.

\subsubsection{Why Hybrid Retrieval Masks the Aggregate Pain Effect}

Semantic similarity (Gemini \texttt{gemini-embedding-001}, 3072-dimensional cosine)
produces consistent and well-calibrated top-10 orderings across the post-incident query
set. For 29 of 31 queries, the semantic score differential between the top-ranked chunk and
its nearest competitor is large enough that the pain multiplier cannot alter the rank order
regardless of the magnitude of the pain differential. The pain term matters only in the
narrow regime where two candidates receive nearly identical semantic similarity
scores---a regime that appeared in exactly 1 of 31 queries evaluated.

This is consistent with findings in the dense retrieval literature: once a high-quality
dense encoder is included in a fused pipeline, sparse signals (including handcrafted boost
signals) have diminishing marginal rank effect \cite{thakur2021beir}. The pain term as
currently implemented is a BM25-tier multiplier applied before RRF fusion; it does not
operate on the post-RRF merged list. A post-RRF re-ranker placement would expose pain to
a different decision boundary and may show larger aggregate effect.

\subsection{Isolated FTS Evaluation and Calibration Test (E10 Follow-up)}
\label{sec:calibration-test}

To isolate pain from semantic dominance, we ran two follow-up ablations.

\textbf{FTS-only ablation ($n$=60 all queries, no semantic layer):}
$\Delta$ nDCG@10 $= +0.0061$, 95\% CI $[-0.0036, +0.0218]$ --- INSIGNIFICANT. Pain
provides no detectable aggregate lift even when the Gemini semantic layer is removed. This
rules out semantic masking as the sole explanation for the hybrid-mode null result.

\textbf{Pain calibration test (4 distributions, E10 follow-up, 2026-05-04).} We
hypothesized that the 89\% of chunks with default $\text{pain}=0.2$ was the limiting
factor. To test this, we constructed three artificial pain distributions and evaluated all
four (real plus three artificial) on the identical FTS+pain pipeline (no Gemini), using
$n=60$ golden queries, 10,000-resample bootstrap, seed=42, on a read-only snapshot
(\texttt{nox-mem-snapshot-20260504-0616.db}). The full $n=60$ set (R01b $\cup$ R01c) is
used here because this is a one-shot ablation post-tuning, not part of the iterative
configuration sprint that froze the salience coefficients; R01c was \emph{not} consulted
during pain-distribution design or coefficient selection. The post-cure $n$=60 R01c-v1.1 corpus
is now the canonical baseline (\S5.1).

\begin{table}[t]
  \caption{Pain calibration test --- pairwise comparisons against real distribution
    ($n$=60 queries, FTS+pain only, bootstrap 95\% CI).}
  \label{tab:calibration}
  \centering
  \resizebox{\textwidth}{!}{%
\begin{tabular}{lllll}
    \toprule
    Distribution & Avg pain & $\Delta$ vs real & 95\% CI & Verdict \\
    \midrule
    real (89\% @ 0.2) & 0.235 & baseline & --- & --- \\
    uniform $[0.1, 1.0]$ & 0.548 & $-0.0148$ & $[-0.041, +0.011]$ & DIRECTIONAL \\
    bimodal $\{0.1, 1.0\}$ & 0.551 & $-0.0087$ & $[-0.038, +0.020]$ & INSIGNIFICANT \\
    log-scale $[0.01, 10.0]$ & 1.447 & $-0.0095$ & $[-0.041, +0.022]$ & INSIGNIFICANT \\
    \bottomrule
  \end{tabular}}
\end{table}

All artificial spreads underperform the real pain distribution. H1, H2, and H3 are all
refuted:

\begin{itemize}
  \item \textbf{H1 [REFUTED].} Real pain (89\% default) is not the limiting factor:
    all three artificial spreads underperform real, with $\Delta$ vs real ranging
    $[-0.0148, -0.0087]$ (Table~\ref{tab:calibration}); none reaches significance,
    confirming calibration spread does not explain the null result.
  \item \textbf{H2 [REFUTED].} Bimodal $\{0.1, 1.0\}$ does not outperform real
    ($\Delta = -0.0087$, INSIGNIFICANT, 95\% CI $[-0.038, +0.020]$). Maximum contrast
    between trivial and outage severities provides no measurable benefit beyond the
    production calibration.
  \item \textbf{H3 [REFUTED].} Log-scale $[0.01, 10.0]$ does not outperform real
    ($\Delta = -0.0095$, INSIGNIFICANT, 95\% CI $[-0.041, +0.022]$). A two-orders-of-
    magnitude wider dynamic range is not the limiting factor.
\end{itemize}

\textbf{Real root cause identified.} Inspection of per-query FTS recall confirms that
55/60 queries (92\%) have \textbf{zero FTS recall on golden chunks}---gold candidates do
not appear in the BM25 candidate pool at all. Pain is a multiplier applied to retrieved
candidates; it cannot affect rank order for queries where the correct chunk is never
retrieved. This is a \textbf{BM25 recall ceiling}, not a calibration limitation.

\textbf{Framing of Contribution~1.} The empirical ablation establishes that the pain
dimension is a \textbf{secondary modulator} effective in tied-semantic regimes
(Q55, $\Delta = +0.349$) rather than a primary ranking signal across all queries. The
design contribution---operationalizing incident severity as a typed schema field and
retrieval multiplier---remains valid. The Q55 case study provides the clearest evidence
that the mechanism functions as intended. The aggregate non-significance reflects the
dominance of the Gemini semantic layer in the hybrid pipeline, not a failure of the pain
construct.

\subsection{Cross-Agent Intelligence Quantification (E12)}

\textbf{Pre-registered hypothesis (E12):} At least 10\% of top-10 retrieved chunks for
any given agent's queries will originate from a different agent's memory namespace,
demonstrating empirically that the shared-canonical design produces cross-agent knowledge
transfer in practice.

\textbf{Storage-level result (confirmed, 2026-05-04).} Direct inspection of the production
database ($n$=61,257 chunks, 2026-05-04) shows that 99.92\% of all chunks are not
partitioned by agent identity.

\begin{table}[t]
  \caption{Cross-agent storage quantification ($n$=61,257 chunks, prod DB, 2026-05-04).}
  \label{tab:crossagent}
  \centering
  \resizebox{\textwidth}{!}{%
\begin{tabular}{lllll}
    \toprule
    Origin class & Chunks & \% & Sharing status \\
    \midrule
    graphify + workspace dumps (\texttt{other}) & 59,772 & 97.58\% & Shared-eligible \\
    docs / specs (\texttt{shared}) & 1,435 & 2.34\% & Shared-eligible \\
    nox agent-private memory & 44 & 0.07\% & Agent-owned \\
    atlas / boris / cipher / forge / lex (combined) & 5 & 0.01\% & Agent-owned \\
    \textbf{Total} & \textbf{61,257} & --- & --- \\
    \textbf{Shared-eligible total} & \textbf{61,207} & \textbf{99.92\%} & --- \\
    \bottomrule
  \end{tabular}}
  \begin{minipage}{\linewidth}
    \vspace{4pt}\footnotesize
    Figures derived from \texttt{SELECT source\_file, COUNT(*) FROM chunks GROUP BY ...} on
    prod READ-ONLY DB (2026-05-04). Agent-private chunks identified by \texttt{source\_file}
    path matching \texttt{memory/agents/<name>/} prefix.
  \end{minipage}
\end{table}

The 99.92\% shared-canonical figure is the single strongest quantitative claim for
Contribution~3. Under the MemGPT/Letta per-agent isolated design \cite{packer2023memgpt},
six agents with comparable memory volumes would maintain six separate corpora with 0\%
sharing---each agent loses access to lessons learned by the other five. nox-mem inverts
this entirely by design.

\textbf{Retrieval-level result: DEFERRED.} The \texttt{search\_telemetry} table does not
include a \texttt{requesting\_agent} column; the schema migration was planned but not
deployed within the W2 window. This migration is documented as a backlog item
(E12-followup: \texttt{ALTER TABLE search\_telemetry ADD COLUMN requesting\_agent TEXT;})
and will enable the retrieval-level quantification after two weeks of telemetry
accumulation.

\begin{table}[t]
  \caption{Cross-agent retrieval attribution matrix ($6{\times}6$). \textsc{[Deferred:
    requires \texttt{search\_telemetry} migration]}}
  \label{tab:attribution}
  \centering
  {\footnotesize
  \setlength{\tabcolsep}{3pt}
\begin{tabular}{lllllll}
    \toprule
    & Orig: nox & Orig: atlas & Orig: boris & Orig: cipher & Orig: forge & Orig: lex \\
    \midrule
    Req: nox    & ---    & [D] & [D] & [D] & [D] & [D] \\
    Req: atlas  & [D] & ---    & [D] & [D] & [D] & [D] \\
    Req: boris  & [D] & [D] & ---    & [D] & [D] & [D] \\
    Req: cipher & [D] & [D] & [D] & ---    & [D] & [D] \\
    Req: forge  & [D] & [D] & [D] & [D] & ---    & [D] \\
    Req: lex    & [D] & [D] & [D] & [D] & [D] & ---    \\
    \bottomrule
  \end{tabular}}
\end{table}

% ---------------------------------------------------------------
% SECTION 6 — DISCUSSION
% ---------------------------------------------------------------
\section{Discussion}

\subsection{What Worked}

Three contributions show empirical or operational validation at the time of writing.

\textbf{Hybrid retrieval pipeline necessity (\S5.1).} The clearest finding in R01c-v1.1 is not
a marginal improvement but a categorical boundary: FTS5 BM25 achieves nDCG@10 $= 0.0000$
(effectively zero) on natural-language queries over the operational corpus. This validates
the hybrid design not as an optimization choice but as an architectural requirement. The
gap of 58.3~pp (absolute)---a 100\% relative reduction in FTS vs.\ hybrid (FTS5 returns zero matches; the absolute $\Delta$ is the entire signal)---confirms the
claim stated in \S3.3: for an operational corpus where queries are issued in natural
language and documents contain domain-specific terminology, hybrid retrieval with a semantic
layer is the minimum viable design.

\textbf{Shadow discipline as incident prevention (\S3.5, \S4.2).} The shadow-mode
architecture prevented at least one class of production regression during the evaluation
period. The incident of 2026-04-25 (\S6.2) involved a ranking-affecting change reaching
production without validation. The subsequent codification of shadow discipline as a
seven-day mandatory gate---enforced via cron and \texttt{/api/health}---means that future
incidents of this class would be detected in shadow telemetry before activation. During the
Fase~1.7b-b salience shadow period, the collected telemetry over 191 promotion candidates,
16,608 review candidates, and 45,743 archive candidates provided the distribution analysis
required for an informed activation decision.

\textbf{Edge typing enum coverage recovery (\S4.4).} Enum coverage rate improved from
14\% to 56\% following the three-path defensive normalization ($4\times$ improvement);
equivalently, the \texttt{unknown} rate decreased from 86\% to 44\% on $n$=100 sampled
relations. This directly enables blast-radius queries
(\texttt{impact~\textless entity\textgreater}) that were practically unusable before the
fix. Note: this is a self-reported coverage rate; the 95\% Wilson CI $[46$--$66\%]$ is a
valid proportion CI on the coverage metric, not a classification-accuracy CI (see \S6.3).

\subsection{What Did Not Work: Incidents That Shaped the Architecture}

\textbf{Incident 2026-04-25: the reindex without dry-run.} At 22:03 on 2026-04-25, a
scheduled end-of-day cron job executed \texttt{nox-mem reindex} without a dry-run flag
against the production database. The reindex routine, using the generic
\texttt{ingestFile()} path rather than the entity-aware \texttt{ingestEntityFile()} router,
processed 183 entity files and stripped their \texttt{section}, \texttt{retention\_days},
and \texttt{section\_boost} annotations---years of structured metadata replaced with default
values in under two minutes. No error was logged. No alert fired. The database obeyed the
instruction correctly. This incident motivated Feature~F02 (the \texttt{withOpAudit()}
wrapper with atomic snapshot), the \texttt{--dry-run} flag on all destructive operations
(A5), and the ingest router (A2) that prevents \texttt{ingestFile()} from processing entity
files without the entity-specific handling path.

\textbf{Incident 2026-05-01: sed on a binary file.} A sweep script applied
\texttt{sed~-i} to a file pattern that inadvertently matched the production SQLite
database. The \texttt{sed} command treated the database as a text file, corrupting page
boundaries across the 1~GB file and eight backup copies. Recovery required a pre-vacuum
backup that had been placed outside the sweep scope for an unrelated reason. This incident
motivated the operational rule: ``never \texttt{sed~-i} on binary files; filter patterns
to \texttt{.json|.md|.sh|.txt|.jsonl|.env} only.'' Both incidents illustrate the paper's
central thesis: the system's architecture was shaped not by theoretical design but by
operational failure. The schema carries their scars.

\subsection{Limitations}

\textbf{Internal-curator bias.} The primary evaluation (R01c-v1.1, $n$=60) was authored by the
same individual who designed and built the system. This is a significant construct validity
risk. We apply four mitigations (\S4.1): the held-out R01c subset, external corpora with
third-party relevance judgments (BEIR TREC-COVID, Stack Exchange), six negative queries
testing specificity, and 10 BEIR TREC-COVID queries evaluated as a cross-curator set (E11,
0\% vocabulary overlap with internal golden set). However, we acknowledge that these
mitigations do not fully eliminate curator bias.

\textbf{Manual pain annotation.} The \texttt{pain} field is currently annotated by hand,
using calibration heuristics described in \S4.3. This introduces two forms of bias: the
annotator may unconsciously assign higher pain to incidents they remember as costly; and
annotation coverage is currently limited to incident-derived entity files.

\textbf{Pain dimension is recall-ceiling-limited.} Empirical ablations (\S5.4) show pain
provides directional but not statistically significant aggregate effect, with one exception
(Q55, $\Delta=+0.349$). Follow-up calibration tests (\S5.5) refute the hypothesis that
distribution spread is the limiting factor. The actual constraint is \textbf{BM25 recall
ceiling}: 92\% of golden queries (55/60) do not surface their gold chunks via lexical
retrieval at all, leaving pain with no candidates to re-rank. Future work should:

\begin{enumerate}
  \item \textbf{Co-optimize semantic recall with pain re-annotation.} Pain re-ranker as a
    post-RRF stage, not a pre-fusion BM25 multiplier.
  \item \textbf{Pain as semantic confidence modulator.} Use pain to break ties when
    semantic cosine score differences between top candidates fall below a threshold
    (e.g., $|\Delta_\text{sem}| < 0.02$).
  \item \textbf{Empirical pain coverage extension.} Expand pain annotation beyond the
    default 0.2 via LLM-driven automatic pain classification (\S6.5).
\end{enumerate}

\textbf{Cross-agent retrieval quantification incomplete.} The storage-level quantification
(99.92\% shared, \S5.6) is confirmed. However, the retrieval-level quantification---the
pre-registered claim that $\geq$10\% of top-10 results cross agent boundaries---cannot be
computed because the \texttt{search\_telemetry} table lacks a \texttt{requesting\_agent}
column. This migration is documented as E12-followup.

\textbf{Short corpus horizon.} The production corpus spans approximately three months
(March--May 2026). A six-month or twelve-month evaluation would provide stronger evidence
for the salience formula's temporal component.

\textbf{Single-author validation.} No inter-rater reliability study was conducted for the
golden query relevance judgments. This means that the nDCG scores cannot be compared
directly with benchmarks that use multi-judge relevance panels.

\textbf{Edge-typing metric is a self-reported enum coverage rate, not a classification
accuracy.} The reported improvement (14\% $\to$ 56\%, $4\times$ gain, $n$=100, 95\%
Wilson CI $[46$--$66\%]$) measures the proportion of new KG relations where the LLM
committed to one of the seven typed values rather than falling back to
\texttt{unknown}. No human annotation file was produced; no inter-rater agreement study
(Cohen's $\kappa$) was conducted. Future work should validate downstream graph quality via
human annotation of a stratified sample (e.g., $n$=100 with two annotators, target
Cohen's $\kappa \geq 0.6$ per Landis and Koch \cite{landis1977measurement}).

\textbf{Pain calibration as engineering choice.} The pain dimension values (0.1, 0.3, 0.5,
0.7, 1.0), the $10\times$ spread, and the multiplicative aggregation form are engineering
choices motivated by operational practice \cite{pagerduty2023severity,beyer2016site}. We
do not claim psychometric or biological validity. The calibration spread ablation
(\S5.5) found that no artificial spread outperforms the real distribution; the binding
constraint is BM25 recall, not calibration. Aggregation form ablation (additive
vs.\ multiplicative) remains as future work.

\subsection{Cost and Compute Profile}

\textbf{Embedding cost.} The operational corpus contains approximately 62K chunks at a
mean of roughly 400 tokens per chunk, yielding an estimated 25M tokens indexed over four
months. At Gemini \texttt{gemini-embedding-001} pricing of \$0.025/1M input tokens
(2026-Q1 \cite{google2026pricing}), total indexing cost is approximately
\textbf{\$0.62 amortized} over the evaluation period. Per-query embedding cost:
$30 \times \$0.025 / 1{,}000{,}000 \approx \$0.00000075$---effectively zero.

\textbf{Total monthly OPEX.} Gemini embedding (queries): $<\$0.05$ + VPS base plan:
approximately \$10/month = \textbf{approximately \$11/month all-in} for a 62K-chunk,
6-agent, always-on production memory system.

\begin{table}[h]
  \caption{Cost comparison: nox-mem vs.\ alternative memory architectures (estimated
    2026-Q1 pricing, small single-team deployment).}
  \label{tab:cost}
  \centering
  \resizebox{\textwidth}{!}{%
\begin{tabular}{llll}
    \toprule
    System & Indexing (one-time) & Per-query & Monthly OPEX \\
    \midrule
    \textbf{nox-mem (this work)} & $\sim$\$0.62 & $\sim$\$0.00000075 & $\sim$\$11 \\
    MemGPT with GPT-4 \cite{packer2023memgpt} & \$0 & \$0.001--\$0.01 (est.) & \$30--\$300 \\
    GraphRAG \cite{edge2024graphrag} & \$10--\$100 per 100K docs & $\sim$\$0.0001 & \$20--\$50 \\
    Mem0 hosted \cite{chhikara2025mem0} & n/a (managed) & bundled & \$20--\$200 \\
    Pure GPT-4o long-context & \$0 & \$0.01--\$0.10/session & \$100--\$1{,}000+ \\
    \bottomrule
  \end{tabular}}
  \begin{minipage}{\linewidth}
    \vspace{4pt}\footnotesize
    MemGPT, GraphRAG, Mem0, and GPT-4o figures are estimated based on documented LLM call
    patterns and do not represent vendor-audited cost disclosures. nox-mem figures are
    measured from production telemetry and API invoices over the 4-month evaluation period.
    Labor cost is the dominant real cost and is excluded; infrastructure OPEX is
    \$11/month.
  \end{minipage}
\end{table}

\subsection{Threats to Validity}

\textbf{Construct validity.} The golden queries (R01b) were designed to reflect operational
retrieval needs---``what was the fix for the gateway crash?'' rather than paper-style
information-need queries. Comparison with external baselines on BEIR (\S5.2) addresses
this partially.

\textbf{External validity.} All internal results (\S5.1--5.2) were collected on a
technology and operations corpus authored by a single software practitioner. The hybrid
pipeline's advantage over FTS-only may not transfer to corpora with different term
distribution properties (e.g., legal documents with precise terminology may show stronger
BM25 performance).

\subsection{Future Work}

\textbf{Automated pain classification.} An LLM-driven incident classifier trained on the
existing pain-annotated chunks as a few-shot signal could extend pain coverage to the full
corpus and reduce annotation bias. This is designated as deferred feature~D02 in the
project roadmap.

\textbf{Cross-encoder reranker (D01).} A cross-encoder reranker applied post-RRF would
likely improve precision on the top-3 results. This feature is gated on
R01c $\geq 0.6$ nDCG@10, following the shadow-mode discipline.

\textbf{Multi-tenant productization (P01).} The shared-canonical design (\S3.7) is not
suitable for multi-tenant SaaS environments. The P01 roadmap item requires a
tenant-isolation layer above the shared corpus, likely via row-level security and per-tenant
\texttt{source\_file} namespacing.

We invite verification and contributions via the public repository at
\url{https://github.com/totobusnello/memoria-nox}, which includes the full code, evaluation
harness, golden query set ($n$=60), and 4-month incident log under MIT license.

% ---------------------------------------------------------------
% SECTION 7 — CONCLUSION
% ---------------------------------------------------------------
\section{Conclusion}

Agent memory is not a retrieval engineering problem. It is an operational discipline
problem. Systems fail silently because ranking changes enter production without validation,
because incident severity is treated as a logging concern rather than a retrieval signal,
and because agents in the same deployment live in context silos that never communicate.
Better embeddings do not fix any of these failure modes. Architecture does.

This paper has described three contributions that address these failure modes directly.
First, \textbf{pain-weighted salience}---$\texttt{salience} = \texttt{recency} \times
\texttt{pain} \times \texttt{importance}$---models incident severity as a first-class
retrieval signal, making a production-outage lesson from six months ago more retrievable
than a minor note updated yesterday. To our knowledge, no prior memory system paper
includes this dimension; the closest related work (GraphRAG, Mem0, MemGPT, A-MEM, HiRAG,
Cognee) models recency and structure but not cost. Second, \textbf{enforced shadow
discipline}---a mandatory seven-day telemetry comparison gate before any ranking-affecting
change reaches production---converts a documentation best practice into an architectural
guarantee. The incident of 2026-04-25 is the counterfactual: a ranking change entered
production without this gate, and 183 entities lost their structured metadata without
alerting. Third, \textbf{shared-canonical multi-agent design} enables cross-agent
knowledge transfer without federation overhead, allowing six agents operating in distinct
domains to benefit from each other's learned context by design.

The empirical evidence supports the hybrid pipeline as a minimum viable requirement
(nDCG@10 $0.5831 \pm 0.0046$ vs $0.0000 \pm 0.0000$ for FTS-only on natural-language
queries, $n$=60 3-run mean post-cure R01c-v1.1; absolute gap 58.3~pp). The BM25 Pyserini comparison is
confirmed: nox-mem hybrid achieves $4.0\times$ the nDCG@10 of the strongest tuned BM25
baseline (+43.6~pp absolute), substantially exceeding the pre-registered threshold. The
shared-canonical storage architecture is confirmed at 99.92\% sharing ($n$=61,302 chunks),
vs.\ 0\% sharing under any per-agent isolated design by construction, achieving production OPEX of approximately
\$11/month all-in for the full 6-agent deployment (\S6.4). The E10 pain ablation (\S5.4)
is executed and reported: the aggregate result is DIRECTIONAL, NOT SIGNIFICANT
($\Delta = +0.0065$, 95\% CI $[-0.0143, +0.0338]$, $n=31$); the Q55 case study provides
positive evidence that pain provides meaningful lift ($\Delta = +0.349$) in the
tied-semantic regime. We characterize pain as a secondary modulator rather than a primary
retrieval signal in hybrid mode. One deferred experiment---the E12 retrieval-level
cross-agent quantification---is documented transparently in \S6.3 and does not alter the
architectural contributions. Note that the current nDCG@10 of $0.5831 < 0.6$, which is below
the threshold at which the optional cross-encoder reranker becomes warranted; this addition
is therefore deferred to future work (\S6.5).

\textbf{Measurement instrument matures alongside the system.} Between paper draft v1.0
(2026-05-04, R01b $n$=50) and v1.1 (2026-05-07, R01c-v1.1 $n$=60), the underlying retrieval
pipeline (FTS5 + Gemini 3072d + RRF $k$=60) was held constant; no code change was made to
search, ranking, or fusion. Yet hybrid nDCG@10 improved from $0.5213 \pm 0.0004$ to
$0.5831 \pm 0.0046$ (+11.9\% relative), Recall@10 from $0.6800$ to $0.7667$ (+12.8\%), and
the absolute $\Delta$ over FTS5 vanilla from $+0.509$ to $+0.583$. The improvement came
\emph{entirely from measurement hygiene plus corpus growth}: 11 queries with empty
\texttt{expected\_chunk\_ids} (genuine documentation gaps where no chunk in the corpus
answers the question) were correctly reclassified to \texttt{category=negative}, where they
score 0.000 in any retrieval system; 3 orphan IDs that referenced chunks deleted by
re-ingest events were corrected; 3 temporal queries (Q70/Q87/Q88) were extended with
newly available timeline events; and the operational corpus grew from 61{,}257 to
61{,}302 chunks with KG \texttt{evidence\_chunk\_id} coverage rising from $0.47\%$ to
$4.92\%$ via incremental nightly extraction.
We treat the gold standard as a versioned research artifact rather than a one-time snapshot,
in keeping with the open-evaluation norm advocated by \cite{rogers2021just} and broader IR
practice of iterative relevance-judgment refinement. The v1.0
numbers are preserved in the \texttt{v1.0.0} git tag for full reproducibility, and the
disclosure (\S5.1, ``Gold standard revision note'') makes both versions available
side-by-side. The implication for IR practitioners: gold-set hygiene---moving doc-gap
queries to negatives, repairing orphan references after corpus rewrites---can deliver
double-digit nDCG@10 gains without touching the model, embedding, or retrieval code. This
is best practice, not artifact, when the methodology is auditable and pre-registered. We
recommend evaluating eval harnesses on the same shadow-discipline timeline (\S3.5) as
ranking changes themselves.

Beyond nox-mem specifically, pain-weighted salience and shadow discipline are
\textbf{transferable concepts}. Any persistent memory system---regardless of implementation
stack---can adopt a severity annotation field and enforce a shadow validation gate before
ranking changes activate. These ideas require no new model, no new architecture, and no
GPU. They require only the discipline to instrument what already exists and the patience to
watch before activating.

We invite verification and contributions via the public repository.

% ---------------------------------------------------------------
% APPENDICES
% ---------------------------------------------------------------
\appendix

\section{Implementation Details}

\subsection{TypeScript Stack}

The system is implemented in TypeScript (Node.js 22, strict mode, ESM modules) with the
following primary dependencies:

\begin{itemize}
  \item \texttt{better-sqlite3} --- synchronous SQLite interface; all DB operations are
    single-file transactions
  \item \texttt{sqlite-vec} --- vector extension for SQLite enabling cosine similarity
    search over 3072-dimensional Gemini embeddings
  \item \texttt{@google/generative-ai} --- Gemini API client for both embedding
    (Gemini embedding-001, 3072-d) and LLM extraction (Gemini 2.5 Flash Lite for agent
    infra; Gemini 2.5 Flash for KG extraction)
  \item \texttt{inotifywait} --- filesystem watch for automatic ingest on file change
\end{itemize}

\textbf{Schema migration history (v1--v12).}

\begin{table}[!ht]
  \centering
  \caption{Schema migration history (v1--v12)}
  \label{tab:schema-migration}
  \resizebox{\textwidth}{!}{%
\begin{tabular}{ll}
    \toprule
    Version & Key change \\
    \midrule
    v1--v3 & Initial chunks + FTS5 design \\
    v4--v5 & \texttt{vec\_chunks} + \texttt{vec\_chunk\_map} (sqlite-vec) \\
    v6--v7 & \texttt{kg\_entities} + \texttt{kg\_relations} \\
    v8     & \texttt{retention\_days} typed retention policy \\
    v9     & \texttt{pain} field (REAL DEFAULT 0.2) \\
    v10    & \texttt{section} + \texttt{section\_boost} for entity file format \\
    v11    & \texttt{search\_telemetry} eval harness (+4 columns) \\
    v12    & \texttt{ops\_audit} table + status enum enforcement triggers \\
    \bottomrule
  \end{tabular}}
\end{table}

All migrations were applied to the production database without downtime.

\subsection{Entry Points}

\begin{itemize}
  \item \textbf{CLI:} \texttt{dist/index.js} (26+ subcommands including \texttt{search},
    \texttt{ingest}, \texttt{ingest-entity}, \texttt{reindex}, \texttt{vectorize},
    \texttt{kg-extract}, \texttt{reflect}, \texttt{crystallize})
  \item \textbf{MCP Server:} 16 tools including \texttt{nox\_mem\_search},
    \texttt{kg\_build}, \texttt{cross\_search}, \texttt{reflect}
  \item \textbf{HTTP API:} port 18802, endpoints
    \texttt{/api/health}, \texttt{/api/search}, \texttt{/api/kg},
    \texttt{/api/kg/path}, \texttt{/api/agents},
    \texttt{/api/cross-kg}, \texttt{/api/reflect},
    \texttt{/api/procedures} + \texttt{POST /api/crystallize}
\end{itemize}

% ------------------------------------------------------------------
\section{Shadow Case Study --- Fase~1.7b-b Salience Activation}

\subsection{Timeline}

\begin{itemize}
  \item \textbf{2026-04-18:} Pain-weighted salience formula implemented;
    \texttt{NOX\_SALIENCE\_MODE=shadow} set in production
  \item \textbf{2026-04-18 to 2026-04-25:} Seven-day shadow telemetry collection
  \item \textbf{2026-04-25:} Distribution analysis: 191 promotion candidates, 16,608
    review candidates, 45,743 archive candidates; distribution shows clear separation
    across tiers
  \item \textbf{2026-04-25:} Decision to advance to \texttt{NOX\_SALIENCE\_MODE=active};
    activation logged in \texttt{ops\_audit}
\end{itemize}

\subsection{Telemetry Summary}

\begin{figure}[!ht]
  \centering
  \includegraphics[width=\textwidth]{figures/figure2-salience-pipeline.pdf}
  \caption{Salience pipeline. Each chunk is annotated at ingest with three
    signals---\texttt{retention\_days}, \texttt{section}, and \textbf{pain}
    (0.1 trivial to 1.0 prod-outage). Salience
    $= \texttt{recency} \times \texttt{pain} \times \texttt{importance}$ is computed and
    exposed via \texttt{/api/health.salience} for at least seven days before activation.}
  \label{fig:salience}
\end{figure}
\FloatBarrier

\begin{table}[h]
  \caption{Telemetry distribution during Fase~1.7b-b shadow period.}
  \centering
  \resizebox{\textwidth}{!}{%
\begin{tabular}{lll}
    \toprule
    Tier & Count & \% of total \\
    \midrule
    Promoted (\texttt{new\_score} $\geq$ threshold) & 191 & 0.30\% \\
    Review (threshold range) & 16,608 & 25.88\% \\
    Archive (\texttt{new\_score} $<$ threshold) & 45,743 & 71.26\% \\
    \textbf{Total shadow observations} & \textbf{64,180+} & --- \\
    \bottomrule
  \end{tabular}}
  \begin{minipage}{0.9\linewidth}
    \vspace{4pt}\footnotesize
    The 3-tier distribution reflects the expected long-tail structure of an operational
    corpus---most chunks are low-salience documentation, a minority are high-salience
    incident lessons.
  \end{minipage}
\end{table}

\subsection{Counterfactual: Incident 2026-04-25}

The incident of 2026-04-25 occurred the same day the shadow period ended---a coincidence
that underscores the motivation for the shadow gate. The reindex operation that damaged 183
entity records is precisely the class of ranking-affecting change that shadow mode is
designed to catch: it altered \texttt{section} and \texttt{section\_boost} annotations,
which feed directly into the salience computation. Had the post-incident reindex run before
the shadow gate was in place, the damage would have been invisible until a user noticed
degraded retrieval quality for entity queries.

\subsection{Activation Decision Rationale}

The distribution analysis showed that the pain-weighted salience formula was doing the
intended work: high-pain chunks (incident lessons, security decisions, production outage
records) consistently landed in the promotion tier, while trivial notes and meeting
summaries landed in the archive tier. The formula was activated because the distribution
matched the design intent, not because the absolute nDCG numbers improved (those were
measured separately in R01b).

% ------------------------------------------------------------------
\section{Reviewer-Friendly Feature Comparison Table}

See Table~\ref{tab:comparison} (\S2.5) for the full seven-axis architectural comparison.
The five-dimension summary below distills the axes most relevant for reviewer quick
reference; scores align with the 5/7 subset of Table~\ref{tab:comparison} that excludes
corpus scale and third-party benchmark coverage.

\begin{table}[h]
  \caption{Detailed architectural feature comparison (reviewer-friendly).
    Five axes shown; the full 7-axis comparison (including corpus scale \textgreater100K
    and 3rd-party benchmark) is in Table~\ref{tab:comparison} of \S2.5. nox-mem covers
    5/7 axes; $\times$ on corpus scale and 3rd-party benchmark.}
  \centering
  \small
  \resizebox{\textwidth}{!}{%
\begin{tabular}{llllll}
    \toprule
    System & KG native & Hybrid retr. & Eval harness & Multi-agent & Shadow disc. \\
    \midrule
    \textbf{nox-mem (this work)} & closed-enum, 7 & FTS5+Gemini+RRF & nDCG/MRR/Recall & shared canonical & $\geq$7d enforced \\
    GraphRAG \cite{edge2024graphrag} & community det. & via KG & No & No & No \\
    MemGPT/Letta \cite{packer2023memgpt} & No & embed.-first & No & per-agent & No \\
    Mem0 \cite{chhikara2025mem0} & Optional (v2) & No --- vector & LOCOMO & user\_id & No \\
    A-MEM \cite{xu2025amem} & Zettelkasten & semantic & No & No & No \\
    HiRAG \cite{huang2025hirag} & hierarchical & multi-level & task-specific & No & No \\
    Cognee \cite{topoteretes2024cognee} & ECL & hybrid & ad-hoc & Optional & No \\
    LangChain Memory & No & key-value & No & session\_id & No \\
    \bottomrule
  \end{tabular}}
\end{table}

% ------------------------------------------------------------------
\section{Incident Case Study --- Formal Write-Up (E13)}

\subsection{Incident Summary}

\textbf{Date:} 2026-04-25, 22:03 BRT \\
\textbf{Severity:} HIGH (data integrity; no user-facing downtime) \\
\textbf{Root cause:} End-of-day cron job executed \texttt{nox-mem reindex} without
dry-run or pre-operation snapshot. The generic \texttt{ingestFile()} path processed 183
entity files, stripping \texttt{section}, \texttt{retention\_days}, and
\texttt{section\_boost} annotations. \\
\textbf{Recovery:} Manual re-ingestion via \texttt{ingest-entity} for all 183 files;
\texttt{withOpAudit()} wrapper retroactively applied as preventive measure. \\
\textbf{Time to detect:} $\approx$18 minutes (via
\texttt{/api/health.sectionDistribution} check). \\
\textbf{Time to recover:} $\approx$45 minutes.

\subsection{Contributing Factors}

\begin{enumerate}
  \item The end-of-day cron script (cron ID \texttt{ee15b430}, 22:00 BRT, step 11)
    invoked \texttt{nox-mem reindex} without the \texttt{--dry-run} flag.
  \item The reindex command used the generic \texttt{ingestFile()} router, which does not
    distinguish entity files from plain markdown.
  \item No pre-operation snapshot was taken; the daily backup (02:00 BRT) had not yet run
    for the day.
  \item No alert was configured for \texttt{section} annotation coverage drop.
\end{enumerate}

\subsection{Changes Implemented}

\begin{table}[!ht]
  \centering
  \caption{Changes implemented post-incident 2026-04-25}
  \label{tab:incident-changes}
  \resizebox{\textwidth}{!}{%
\begin{tabular}{lll}
    \toprule
    Change & Code artifact & Description \\
    \midrule
    F02 op-audit & \texttt{src/lib/op-audit.ts} & \texttt{withOpAudit()}: VACUUM INTO atomic snapshot \\
    A2 ingest router & \texttt{src/lib/ingest-router.ts} & \texttt{routeIngest()} dispatches entity files correctly \\
    A5 dry-run & \texttt{reindex.ts}, \texttt{consolidate.ts} & \texttt{--dry-run} produces JSON preview without DB mutation \\
    Schema invariant canary & \texttt{check-schema-invariants.sh} & Cron \texttt{*/15min}; Discord alert on deviation \\
    Cron patch & \texttt{ee15b430} step 11 & Replaced \texttt{reindex} with \texttt{consolidate} (entity-aware) \\
    \bottomrule
  \end{tabular}}
\end{table}

\subsection{Lessons Formalized in Architecture}

The \texttt{withOpAudit()} module encodes the lesson that destructive operations must
create a point-in-time snapshot before execution and log their outcome to an append-only
audit table (\texttt{ops\_audit}). The audit table's \texttt{status} field is validated
by DB triggers against a closed enum (\texttt{started}, \texttt{success},
\texttt{failed}, \texttt{crashed}); the triggers block DELETE and UPDATE on rows with
terminal status. Recovery via \texttt{safeRestore()} validates \texttt{user\_version}
match before restoring---a safeguard motivated by a separate incident where a stale WAL
file caused silent corruption on a naive \texttt{cp} restore.

This incident, and the five others documented in \texttt{docs/INCIDENTS.md}, are the
operational substrate from which the paper's architectural contributions were extracted.
They are included here not to confess failure but to demonstrate that the contributions
are grounded in production evidence rather than synthetic benchmark design.


% ---------------------------------------------------------------
% BIBLIOGRAPHY
% ---------------------------------------------------------------
\bibliographystyle{unsrt}
\bibliography{../refs}

% [ORIGINAL MONOLITH PRESERVED BELOW FOR REFERENCE — remove before submission]
\iffalse
% ---------------------------------------------------------------
% SECTION 4 — METHODS (MONOLITH)
% ---------------------------------------------------------------
\section{Methods}

We describe the evaluation framework, shadow-mode methodology, and calibration procedures
used to validate the three primary contributions. All experiments were pre-registered in
\texttt{03-experiments-needed.md} before results were collected, following the
open-evaluation norm advocated by \cite{rogers2021just}.

\subsection{Evaluation Harness}

Our primary evaluation uses \textbf{nDCG@10}, \textbf{MRR} (Mean Reciprocal Rank),
\textbf{Recall@10}, and \textbf{Precision@5} computed over a set of 50 internally curated
golden queries (dataset R01b, fully cured as of 2026-05-03). These metrics follow the
standard IR evaluation methodology described in \cite{manning2008introduction}.

\textbf{Golden query construction.} Queries span eight categories reflecting the
operational nature of the corpus: \texttt{entity} (specific named entities---agents, tools,
decisions), \texttt{procedure} (how-to operational steps), \texttt{concept} (abstract
architectural notions), \texttt{security} (vulnerability and mitigation queries),
\texttt{decision} (architectural choices and their rationale), \texttt{cross-agent}
(questions whose answer originates from a different agent's memory space),
\texttt{temporal} (time-anchored recall, e.g., ``what changed in late April''), and
\texttt{negative} (6 queries, 12\% of set, for which the correct answer is that no
relevant chunk exists---testing specificity against hallucination risk). Each query was
authored by the single curator with a relevance label set ($0 =$ not relevant,
$1 =$ partially relevant, $2 =$ highly relevant) over the top-20 retrieved candidates.

\textbf{Held-out subset (R01c).} Ten queries from R01b are designated held-out: they were
locked before any retrieval tuning and are evaluated only once per major system revision,
functioning as a proxy for external-curator independence. Performance on R01c is reported
separately from the 40-query main set to avoid optimistic bias from iterative query
refinement.

\textbf{Internal-curator bias mitigation.} We acknowledge that golden queries authored by
the same individual who built the system introduce construct validity risk. Three
mitigations are applied: (i)~the held-out R01c subset was frozen before the final tuning
sprint; (ii)~external corpora (BEIR TREC-COVID, Stack Exchange---\S5.3) use third-party
curated relevance judgments; and (iii)~six negative queries test the boundary condition
most susceptible to self-serving bias.

\subsection{Shadow-Mode Methodology}

Any change that affects retrieval ranking in the production system is subject to mandatory
shadow validation before activation. The protocol is enforced architecturally via the
environment variable \texttt{NOX\_SALIENCE\_MODE}, which accepts three values:
\texttt{shadow} (collect both old and new scores in \texttt{search\_telemetry} without
applying the new ranking), \texttt{active} (apply new ranking), and \texttt{off} (disable
the feature entirely).

\textbf{Telemetry collection.} In shadow mode, every search call writes a row to
\texttt{search\_telemetry} containing: \texttt{query\_text} (opt-in,
\texttt{NOX\_SEARCH\_LOG\_TEXT=1}), \texttt{old\_score}, \texttt{new\_score},
\texttt{top\_chunk\_ids}, and \texttt{top\_scores}. This enables offline comparison of
old and new score distributions without exposing users to the changed ranking.

\textbf{Activation gate.} Shadow validation runs for a minimum of seven calendar days.
After the shadow period, the stored distribution is analyzed: if the new score distribution
shows statistically meaningful separation from the old distribution (inspected via
percentile comparison and visual histogram), and if no ranking inversion is detected on a
manually reviewed 10-query spot check, the feature advances to \texttt{active}. The
seven-day minimum is not a guideline---it is a hard constraint codified in the cron
configuration that governs feature activation. This design choice is motivated by the
incident of 2026-04-25, where a ranking-affecting change reached production without any
offline validation period, causing 183 entity records to lose their structured metadata
without triggering any alert (\S6.2).

\textbf{Case study: Fase~1.7b-b salience activation.} During the seven-day shadow period
for the pain-weighted salience formula, the system collected telemetry over 191 promotion
candidates, 16,608 review candidates, and 45,743 archive candidates. The distribution
separated clearly across all three tiers. Only after this distribution analysis did we
advance \texttt{NOX\_SALIENCE\_MODE} from \texttt{shadow} to \texttt{active}. This case
study is documented in detail in Appendix~B.

\subsection{Pain Weighting Calibration}

The \texttt{pain} field is a real-valued annotation in $[0.1,\,1.0]$ attached to each
chunk at ingest time. Annotation is currently manual, using the
\texttt{<!--~pain:~X.X~-->} comment syntax in entity files, and defaults to 0.2 for
unannotated content.

\textbf{Calibration heuristics.} Based on four months of operational experience, the
following calibration anchors were established: 0.1 (trivial notes, meeting summaries with
no operational consequence); 0.2 (default, documentation and informational content);
0.3--0.4 (decisions with moderate reversibility risk); 0.5--0.7 (production incidents with
bounded impact---recoverable within one session); 0.8--0.9 (incidents causing data loss or
multi-hour outages); 1.0 (catastrophic incidents---unrecoverable data loss, multi-day
downtime, or security breach). The calibration is conservative: in ambiguous cases,
annotators are instructed to use the lower bound of the relevant range, then escalate only
if post-incident analysis reveals higher severity.

\textbf{Pain dimension used in the salience formula.} See Equation~\ref{eq:salience}.
where $\texttt{recency} \in [0,1]$ is an exponential decay over \texttt{last\_seen}
timestamp, $\texttt{pain} \in [0.1,1.0]$ is the manual annotation, and
$\texttt{importance} \in [0,1]$ is derived from \texttt{mention\_count} and
\texttt{entity\_type} prior. The multiplicative structure means that a high-pain chunk
remains salient even as its recency decays---which is the core behavioral claim of
Contribution~1.

\textbf{Annotation coverage.} The corpus contains 64,180+ chunks; pain annotation is
currently applied selectively to chunks derived from incident entity files. Future work
includes LLM-driven automatic pain classification over the full corpus (\S6.5).

\subsection{Edge Typing Extraction}

\begin{figure}[t]
  \centering
  \includegraphics[width=\linewidth]{figures/figure4-kg-edge-typing.pdf}
  \caption{KG edge typing flow. Gemini emits free-text relation labels; a defensive
    three-path normalizer collapses them into a closed enum of seven canonical reasons
    enforced by a schema-v12 \texttt{CHECK} constraint. Before fix: 14\% classified,
    86\% unknown. After fix: 56\% classified, 4$\times$ improvement (n=100).}
  \label{fig:edgetyping}
\end{figure}

The knowledge graph relation schema uses a closed-enum field \texttt{relation\_reason}
with seven values: \texttt{depends\_on}, \texttt{derived\_from}, \texttt{opposes},
\texttt{extends}, \texttt{replaces}, \texttt{mentions}, and \texttt{unknown}. The goal of
edge typing is to enable blast-radius queries (e.g., ``what does component X depend on?'')
that are impossible with untyped relations.

\textbf{Prompt design and the unknown-rate problem.} An initial prompt that marked
\texttt{relation\_reason} as an optional field with the instruction ``use \texttt{unknown}
if unsure'' produced 86\% unknown-typed relations across n=100 sampled extractions,
rendering the typed KG practically useless for blast-radius queries. The fix applied a
three-path defensive normalization strategy: (i)~a revised prompt that provides explicit
examples for each of the six non-unknown categories and makes \texttt{unknown} a last
resort rather than a default; (ii)~a code-side defensive map
(\texttt{RELATION\_TYPE\_TO\_REASON}, 24 entries) that normalizes LLM-produced free-text
variants---including PT-BR and EN aliases---to the 7 canonical enum values; and (iii)~a
post-extraction validation pass that re-prompts any row where the LLM output did not match
the closed enum. After this fix, the classification rate improved from 14\% to 56\% on
n=100 sampled relations (4$\times$ improvement); equivalently, the \texttt{unknown} rate
decreased from 86\% to 44\%.

\textbf{Current KG state (2026-05-03).} The production graph contains approximately 402
entities and 544 relations, extracted incrementally by a nightly Gemini 2.5 Flash job.

\subsection{Statistical Methodology}

All retrieval experiments report \textbf{3-run mean $\pm$ standard deviation} with Bessel
correction. The system is operationally deterministic for identical queries against a
static corpus (no stochastic ranking), so run-to-run variance arises primarily from corpus
index state and warm-cache effects; empirically, standard deviation across runs is
consistently below 0.001 for nDCG@10.

For small-N validations---specifically the pain dimension experiment (E10,
n=10--15 post-incident queries)---we additionally report \textbf{bootstrap 95\% confidence
intervals} computed with 10,000 resamples. Given the small sample, bootstrap CI is the
appropriate uncertainty quantification; we do not use asymptotic normal approximations for
$n < 30$.

Effect sizes for ablation experiments (\S5.4) are reported as absolute $\Delta$~nDCG@10
rather than relative percentages, following the recommendation of \cite{fuhr2018some} to
avoid inflating small absolute differences through percentage framing.

% ---------------------------------------------------------------
% SECTION 5 — EXPERIMENTS AND RESULTS
% ---------------------------------------------------------------
\section{Experiments and Results}

We report results across five experimental questions: (5.1) internal corpus baseline
establishing hybrid pipeline necessity; (5.2) comparison against strong external baselines;
(5.3) generalization to external corpora; (5.4) ablation studies isolating each
architectural layer; and (5.5--5.6) targeted validation of the two novel contributions---pain
weighting and cross-agent intelligence. All pre-registered hypotheses from
\texttt{03-experiments-needed.md} are stated before results; pending experiments are marked
\textsc{[Pending: W2]} or \textsc{[Pending: W3]}.

\subsection{Internal Corpus Baseline (R01a/b/c)}

\textbf{Pre-registered hypothesis (R01a):} The hybrid pipeline will outperform FTS-only
BM25 by a substantial margin on natural-language queries over the operational corpus.

\textbf{Result (confirmed, R01b/R01c, 2026-05-03):} Table~\ref{tab:hybrid} shows the
primary comparison. FTS5 vanilla BM25 achieves nDCG@10 $= 0.0123$ (effectively zero) on
natural-language queries against the operational corpus. This is not an artifact of query
phrasing: FTS5 applies AND-strict matching by default, which means any multi-word
natural-language query that does not appear verbatim in the corpus returns zero results.
This is a structural property of the retrieval system, not a tuning failure, and it
confirms that hybrid retrieval is a minimum viable requirement rather than an optimization
for this corpus type.

\begin{table}[t]
  \caption{Hybrid vs.\ FTS-only on internal corpus
    (R01b, n=50 golden queries, 3-run mean $\pm$ std).}
  \label{tab:hybrid}
  \centering
  \begin{tabular}{lllll}
    \toprule
    Approach & nDCG@10 & MRR & Recall@10 & Precision@5 \\
    \midrule
    FTS5 vanilla (BM25) & $0.0123 \pm 0.0000$ & $0.0200 \pm 0.0000$ & $0.0100 \pm 0.0000$ & $0.0040 \pm 0.0000$ \\
    \textbf{nox-mem hybrid (FTS+Gemini+RRF)} & $\mathbf{0.5213 \pm 0.0004}$ & $0.4889 \pm 0.0028$ & $0.6800 \pm 0.0047$ & $0.2640 \pm 0.0000$ \\
    $\Delta$ (hybrid $-$ FTS) & \textbf{+50.9 pp} & --- & --- & --- \\
    \bottomrule
  \end{tabular}
  \begin{minipage}{\linewidth}
    \vspace{4pt}
    \footnotesize
    3-run mean $\pm$ std (Runs \#10/\#11/\#12 for hybrid, \#13/\#14/\#15 for FTS).
    FTS-only near-zero result is structural (AND-strict matching), not a failure of
    parameterization. Hybrid latency: 119.7\,s / 50 queries ($\approx$2.4\,s/query).
  \end{minipage}
\end{table}

\begin{table}[t]
  \caption{Three-run replication stability
    (nox-mem hybrid, R01b/R01c, n=50 per run).}
  \label{tab:replication}
  \centering
  \begin{tabular}{lll}
    \toprule
    Run & nDCG@10 & Notes \\
    \midrule
    Run \#10 & 0.5213 & Stable post-R01b configuration \\
    Run \#11 & 0.5213 & Replication run \\
    Run \#12 & 0.5213 & Replication run \\
    \textbf{Mean $\pm$ std} & $\mathbf{0.5213 \pm 0.0004}$ & Bessel-corrected 3-run mean \\
    \bottomrule
  \end{tabular}
  \begin{minipage}{\linewidth}
    \vspace{4pt}
    \footnotesize
    Runs \#10--\#12 conducted on the stable post-R01b configuration with no ranking changes
    between runs, isolating API-level variance. Earlier diagnostic runs (Run~\#6: 0.714,
    Run~\#7: 0.674) reflected intermediate config changes and are excluded from the headline
    claim. std=0.0004 (0.08\% relative) confirms the system is operationally deterministic
    on a static corpus.
  \end{minipage}
\end{table}

\begin{table}[t]
  \caption{R01b nDCG@10 breakdown by query category (n=50, hybrid, 3-run mean).}
  \label{tab:category}
  \centering
  \begin{tabular}{llll}
    \toprule
    Category & n & nDCG@10 & Notes \\
    \midrule
    entity      & TBD & TBD & \textsc{[Pending: per-category breakdown from R01b]} \\
    procedure   & TBD & TBD & \textsc{[Pending]} \\
    concept     & TBD & TBD & \textsc{[Pending]} \\
    security    & TBD & TBD & \textsc{[Pending]} \\
    decision    & TBD & TBD & \textsc{[Pending]} \\
    cross-agent & TBD & TBD & \textsc{[Pending]} \\
    temporal    & TBD & TBD & \textsc{[Pending]} \\
    negative    & 6   & TBD & Specificity test---correct answer: no relevant chunk \\
    \textbf{All} & \textbf{50} & $\mathbf{0.5213 \pm 0.0004}$ & 3-run mean (Runs \#10/\#11/\#12) \\
    \bottomrule
  \end{tabular}
\end{table}

\subsection{Comparison Against Strong External Baselines (E1+E2+E3)}

\textbf{Pre-registered hypothesis (E1--E3):} The nox-mem hybrid pipeline will maintain a
nDCG@10 advantage of $\geq$10 percentage points over BGE-M3 dense retrieval on the
internal operational corpus (Corpus~A).

This hypothesis is pre-registered prior to collecting results, in accordance with the
open-evaluation norm. The choice of BGE-M3 as the primary comparison point reflects its
status as a strong open-source dense encoder on the MTEB leaderboard
\cite{muennighoff2022mteb}.

\begin{table}[t]
  \caption{System comparison on internal corpus, Corpus~A
    (n=50 golden queries, 3-run mean $\pm$ std). \textsc{[Pending: W2]}}
  \label{tab:external}
  \centering
  \begin{tabular}{lllll}
    \toprule
    System & nDCG@10 & MRR & Recall@10 & Precision@5 \\
    \midrule
    BM25 (Pyserini) \cite{yang2018anserini}              & {[}P{]} & {[}P{]} & {[}P{]} & {[}P{]} \\
    BGE-M3 \cite{chen2024bge}                            & {[}P{]} & {[}P{]} & {[}P{]} & {[}P{]} \\
    E5-mistral-7b-instruct \cite{wang2023improving}      & {[}P{]} & {[}P{]} & {[}P{]} & {[}P{]} \\
    \textbf{nox-mem hybrid (this work)}                  & $\mathbf{0.5213 \pm 0.0004}$ & $0.4889 \pm 0.0028$ & $0.6800 \pm 0.0047$ & $0.2640 \pm 0.0000$ \\
    \bottomrule
  \end{tabular}
\end{table}

\begin{table}[t]
  \caption{Pre-registered directional hypothesis summary (E1+E2, Corpus~A).}
  \label{tab:hypotheses}
  \centering
  \begin{tabular}{lll}
    \toprule
    Comparison & Pre-registered $\Delta$ nDCG@10 & Result \\
    \midrule
    nox-mem vs BGE-M3           & $\geq +10$ pp & \textsc{[Pending: W2]} \\
    nox-mem vs BM25 (Pyserini)  & $\geq +30$ pp (given 0.0123 FTS5 baseline) & \textsc{[Pending: W2]} \\
    nox-mem vs E5-mistral       & $\geq 0$ pp & \textsc{[Pending: W2]} \\
    \bottomrule
  \end{tabular}
\end{table}

\subsection{Cross-Corpus Generalization (E4+E5)}

\textbf{Pre-registered hypothesis (E4+E5):} The nox-mem hybrid architecture will show
positive nDCG@10 on external corpora (BEIR TREC-COVID, Stack Exchange), confirming that
the three-layer pipeline is not overfit to the internal operational corpus.

\begin{table}[t]
  \caption{Cross-corpus nDCG@10 --- BEIR TREC-COVID subset
    (171K chunks, standard 50 BEIR queries). \textsc{[Pending: W3]}}
  \label{tab:beir}
  \centering
  \begin{tabular}{llll}
    \toprule
    System & nDCG@10 & MRR & Recall@10 \\
    \midrule
    BM25 (Pyserini)  & {[}P{]} & {[}P{]} & {[}P{]} \\
    BGE-M3           & {[}P{]} & {[}P{]} & {[}P{]} \\
    E5-mistral-7b    & {[}P{]} & {[}P{]} & {[}P{]} \\
    \textbf{nox-mem hybrid} & {[}P{]} & {[}P{]} & {[}P{]} \\
    \bottomrule
  \end{tabular}
\end{table}

\begin{table}[t]
  \caption{Cross-corpus nDCG@10 --- Stack Exchange 10K subset
    (mixed factoid/how-to/opinion, 50 curated queries). \textsc{[Pending: W3]}}
  \label{tab:stackex}
  \centering
  \begin{tabular}{llll}
    \toprule
    System & nDCG@10 & MRR & Recall@10 \\
    \midrule
    BM25 (Pyserini)  & {[}P{]} & {[}P{]} & {[}P{]} \\
    BGE-M3           & {[}P{]} & {[}P{]} & {[}P{]} \\
    E5-mistral-7b    & {[}P{]} & {[}P{]} & {[}P{]} \\
    \textbf{nox-mem hybrid} & {[}P{]} & {[}P{]} & {[}P{]} \\
    \bottomrule
  \end{tabular}
\end{table}

\subsection{Ablation Studies (E6--E9)}

\textbf{Pre-registered hypothesis (E6--E9):} Each of the four architectural layers (FTS5
lexical, Gemini semantic embeddings, RRF fusion, section boost) contributes positively to
nDCG@10, with each layer's removal causing a $\Delta \geq 0.03$ decrease.

\begin{table}[t]
  \caption{Ablation study on internal corpus, Corpus~A
    (n=50 golden queries, 3-run mean $\pm$ std). \textsc{[Pending: W3]}}
  \label{tab:ablation}
  \centering
  \begin{tabular}{llll}
    \toprule
    Configuration & nDCG@10 & $\Delta$ vs full hybrid & MRR \\
    \midrule
    Full hybrid (baseline)                             & $0.5213 \pm 0.0004$ & --- & $0.4889 \pm 0.0028$ \\
    FTS-only (no semantic, no RRF)                     & $0.0123 \pm 0.0000$ & $-0.509$ & $0.0200 \pm 0.0000$ \\
    FTS + semantic, no RRF (score concat)              & {[}P{]} & {[}P{]} & {[}P{]} \\
    Hybrid, no salience boost                          & {[}P{]} & {[}P{]} & {[}P{]} \\
    Hybrid, no \texttt{section\_boost}                 & {[}P{]} & {[}P{]} & {[}P{]} \\
    \bottomrule
  \end{tabular}
  \begin{minipage}{\linewidth}
    \vspace{4pt}
    \footnotesize
    FTS-only is confirmed from R01b (Table~\ref{tab:hybrid}). The remaining three ablations
    (E7--E9) are pending implementation and execution. They are controlled via environment
    flags: \texttt{NOX\_RRF\_DISABLE=1}, \texttt{NOX\_SALIENCE\_MODE=off},
    \texttt{NOX\_SECTION\_BOOST\_MODE=off}.
  \end{minipage}
\end{table}

\subsection{Pain Dimension Validation (E10)}

\textbf{Pre-registered hypothesis (E10):} On a subset of post-incident golden queries
(queries where the ground-truth answer is a chunk describing a production incident or
costly operational lesson), pain-aware retrieval (current default) will outperform
pain-uniform retrieval ($\texttt{pain}=1.0$ for all chunks) by $\Delta~\text{nDCG@10}
\geq 0.05$.

The pain-uniform counterfactual collapses all chunks to the same pain weight, effectively
reducing the salience formula to
$\texttt{salience} = \texttt{recency} \times \texttt{importance}$. This tests whether the
pain dimension adds independent retrieval signal beyond recency and importance.

\begin{table}[t]
  \caption{Pain-aware vs.\ pain-uniform on post-incident queries
    (n=10--15 queries from R01b, bootstrap 95\% CI). \textsc{[Pending: W3]}}
  \label{tab:pain}
  \centering
  \begin{tabular}{llll}
    \toprule
    Configuration & nDCG@10 & Bootstrap 95\% CI & MRR \\
    \midrule
    Pain-aware (default, $\texttt{pain} \in [0.1, 1.0]$) & {[}P{]} & {[}P{]} & {[}P{]} \\
    Pain-uniform (counterfactual, $\texttt{pain}=1.0$)    & {[}P{]} & {[}P{]} & {[}P{]} \\
    $\Delta$ (pain-aware $-$ pain-uniform)                & {[}P{]} & {[}P{]} & {[}P{]} \\
    \bottomrule
  \end{tabular}
  \begin{minipage}{\linewidth}
    \vspace{4pt}
    \footnotesize
    Bootstrap CI uses 10,000 resamples. Given the small-N design (n=10--15), we report
    bootstrap CI rather than asymptotic normal approximations. The pre-registered threshold
    is $\Delta \geq 0.05$; results below this threshold would require reframing of
    Contribution~1 claims.
  \end{minipage}
\end{table}

\subsection{Cross-Agent Intelligence Quantification (E12)}

\textbf{Pre-registered hypothesis (E12):} At least 10\% of top-10 retrieved chunks for
any given agent's queries will originate from a different agent's memory namespace,
demonstrating empirically that the shared-canonical design produces cross-agent knowledge
transfer in practice.

\begin{table}[t]
  \caption{Cross-agent hit rates
    (\% of top-10 results originating from a different agent's namespace).
    \textsc{[Pending: W3]}}
  \label{tab:crossagent}
  \centering
  \begin{tabular}{llll}
    \toprule
    Requesting agent & Same-agent hits (\%) & Cross-agent hits (\%) & n queries \\
    \midrule
    nox    & {[}P{]} & {[}P{]} & {[}P{]} \\
    atlas  & {[}P{]} & {[}P{]} & {[}P{]} \\
    boris  & {[}P{]} & {[}P{]} & {[}P{]} \\
    cipher & {[}P{]} & {[}P{]} & {[}P{]} \\
    forge  & {[}P{]} & {[}P{]} & {[}P{]} \\
    lex    & {[}P{]} & {[}P{]} & {[}P{]} \\
    \textbf{All agents} & {[}P{]} & {[}P{]} & {[}P{]} \\
    \bottomrule
  \end{tabular}
\end{table}

\begin{table}[t]
  \caption{Cross-agent attribution matrix (6$\times$6, requesting agent $\times$ origin
    agent, \% of total cross-agent hits). \textsc{[Pending: W3]}}
  \label{tab:attribution}
  \centering
  \begin{tabular}{lllllll}
    \toprule
           & Orig: nox & Orig: atlas & Orig: boris & Orig: cipher & Orig: forge & Orig: lex \\
    \midrule
    Req: nox    & ---    & {[}P{]} & {[}P{]} & {[}P{]} & {[}P{]} & {[}P{]} \\
    Req: atlas  & {[}P{]} & ---    & {[}P{]} & {[}P{]} & {[}P{]} & {[}P{]} \\
    Req: boris  & {[}P{]} & {[}P{]} & ---    & {[}P{]} & {[}P{]} & {[}P{]} \\
    Req: cipher & {[}P{]} & {[}P{]} & {[}P{]} & ---    & {[}P{]} & {[}P{]} \\
    Req: forge  & {[}P{]} & {[}P{]} & {[}P{]} & {[}P{]} & ---    & {[}P{]} \\
    Req: lex    & {[}P{]} & {[}P{]} & {[}P{]} & {[}P{]} & {[}P{]} & ---    \\
    \bottomrule
  \end{tabular}
\end{table}

% ---------------------------------------------------------------
% SECTION 6 — DISCUSSION
% ---------------------------------------------------------------
\section{Discussion}

\subsection{What Worked}

Three contributions show empirical or operational validation at the time of writing.

\textbf{Hybrid retrieval pipeline necessity (\S5.1).} The clearest finding in R01b is not
a marginal improvement but a categorical boundary: FTS5 BM25 achieves nDCG@10 $= 0.0123$
(effectively zero) on natural-language queries over the operational corpus. This validates
the hybrid design not as an optimization choice but as an architectural requirement. The
gap of 50.9~pp (absolute)---a 97.6\% relative reduction in FTS vs.\ hybrid---confirms the
claim stated in \S3.3: for an operational corpus where queries are issued in natural
language and documents contain domain-specific terminology that does not match query terms
lexically, hybrid retrieval with a semantic layer is the minimum viable design.

\textbf{Shadow discipline as incident prevention (\S3.5, \S4.2).} The shadow-mode
architecture prevented at least one class of production regression during the evaluation
period. The incident of 2026-04-25 (\S6.2) involved a ranking-affecting change reaching
production without validation. The subsequent codification of shadow discipline as a
seven-day mandatory gate---enforced via cron and \texttt{/api/health}---means that future
incidents of this class would be detected in shadow telemetry before activation.

\textbf{Edge typing recall recovery (\S4.4).} Classification rate improved from 14\% to
56\% following the three-path defensive normalization (4$\times$ improvement);
equivalently, the \texttt{unknown} rate decreased from 86\% to 44\% on n=100 sampled
relations. This directly enables blast-radius queries (\texttt{impact~\textless entity\textgreater}) that were practically unusable before the fix.

\subsection{What Did Not Work: Incidents That Shaped the Architecture}

\textbf{Incident 2026-04-25: the reindex without dry-run.} At 22:03 on 2026-04-25, a
scheduled end-of-day cron job executed \texttt{nox-mem reindex} without a dry-run flag
against the production database. The reindex routine, using the generic
\texttt{ingestFile()} path rather than the entity-aware \texttt{ingestEntityFile()} router,
processed 183 entity files and stripped their \texttt{section}, \texttt{retention\_days},
and \texttt{section\_boost} annotations---years of structured metadata replaced with default
values in under two minutes. No error was logged. No alert fired. The database obeyed the
instruction correctly. This incident motivated Feature~F02 (the \texttt{withOpAudit()}
wrapper with atomic snapshot), the \texttt{-{}-dry-run} flag on all destructive operations
(A5), and the ingest router (A2) that prevents \texttt{ingestFile()} from processing entity
files without the entity-specific handling path.

\textbf{Incident 2026-05-01: sed on a binary file.} A sweep script applied
\texttt{sed~-i} to a file pattern that inadvertently matched the production SQLite
database. The \texttt{sed} command treated the database as a text file, corrupting page
boundaries across the 1~GB file and eight backup copies. Recovery required a pre-vacuum
backup that had been placed outside the sweep scope for an unrelated reason. This incident
motivated the operational rule: ``never \texttt{sed~-i} on binary files; filter patterns
to \texttt{.json|.md|.sh|.txt|.jsonl|.env} only.'' Both incidents illustrate the paper's
central thesis: the system's architecture was shaped not by theoretical design but by
operational failure. The schema carries their scars.

\subsection{Limitations}

\textbf{Internal-curator bias.} The primary evaluation (R01b, n=50) was authored by the
same individual who designed and built the system. This is a significant construct validity
risk. We apply three mitigations (\S4.1): the held-out R01c subset, external corpora with
third-party relevance judgments (BEIR TREC-COVID, Stack Exchange), and six negative queries
testing specificity. However, we acknowledge that these mitigations do not fully eliminate
curator bias; results on external corpora (\S5.3) are the most important check.

\textbf{Manual pain annotation.} The \texttt{pain} field is currently annotated by hand,
using calibration heuristics described in \S4.3. This introduces two forms of bias: the
annotator (the system author) may unconsciously assign higher pain to incidents they
remember as costly, even when the actual retrieval impact is low; and the annotation
coverage is currently limited to incident-derived entity files.

\textbf{Short corpus horizon.} The production corpus spans approximately four months
(March--May 2026). This is sufficient to validate hybrid retrieval and edge typing, but
may underestimate the long-term recall decay problem that pain weighting is designed to
address. A six-month or twelve-month evaluation would provide stronger evidence for the
salience formula's temporal component.

\textbf{Single-author validation.} No inter-rater reliability study was conducted for the
golden query relevance judgments. This means that the nDCG scores cannot be compared
directly with benchmarks that use multi-judge relevance panels.

\subsection{Threats to Validity}

\textbf{Construct validity.} The golden queries (R01b) were designed to reflect
operational retrieval needs---``what was the fix for the gateway crash?'' rather than
paper-style information-need queries. Comparison with external baselines on BEIR (\S5.3)
addresses this partially.

\textbf{External validity.} All internal results (\S5.1--5.2) were collected on a
technology and operations corpus authored by a single software practitioner. The hybrid
pipeline's advantage over FTS-only may not transfer to corpora with different term
distribution properties (e.g., legal documents with precise terminology may show stronger
BM25 performance).

\subsection{Future Work}

\textbf{Automated pain classification.} An LLM-driven incident classifier trained on the
existing pain-annotated chunks as a few-shot signal could extend pain coverage to the full
corpus and reduce annotation bias. This is designated as deferred feature~D02 in the
project roadmap.

\textbf{Cross-encoder reranker (D01).} A cross-encoder reranker applied post-RRF would
likely improve precision on the top-3 results, where the current pipeline's RRF fusion
sometimes ranks partially relevant chunks above highly relevant ones. This feature is gated
on R01c $\geq 0.6$ nDCG@10, following the shadow-mode discipline.

\textbf{Multi-tenant productization (P01).} The shared-canonical design (\S3.6) is not
suitable for multi-tenant SaaS environments. The P01 roadmap item requires a
tenant-isolation layer above the shared corpus, likely via row-level security and per-tenant
\texttt{source\_file} namespacing.

% ---------------------------------------------------------------
% SECTION 7 — CONCLUSION
% ---------------------------------------------------------------
\section{Conclusion}

Agent memory is not a retrieval engineering problem. It is an operational discipline
problem. Systems fail silently because ranking changes enter production without validation,
because incident severity is treated as a logging concern rather than a retrieval signal,
and because agents in the same deployment live in context silos that never communicate.
Better embeddings do not fix any of these failure modes. Architecture does.

This paper has described three contributions that address these failure modes directly.
First, \textbf{pain-weighted salience}---$\texttt{salience} = \texttt{recency} \times
\texttt{pain} \times \texttt{importance}$---models incident severity as a first-class
retrieval signal, making a production-outage lesson from six months ago more retrievable
than a minor note updated yesterday. To our knowledge, no prior memory system paper
includes this dimension; the closest related work (GraphRAG, Mem0, MemGPT, A-MEM, HiRAG,
Cognee) models recency and structure but not cost. Second, \textbf{enforced shadow
discipline}---a mandatory seven-day telemetry comparison gate before any ranking-affecting
change reaches production---converts a documentation best practice into an architectural
guarantee. The incident of 2026-04-25 is the counterfactual: a ranking change entered
production without this gate, and 183 entities lost their structured metadata without
alerting. Third, \textbf{shared-canonical multi-agent design} enables cross-agent
knowledge transfer without federation overhead, allowing six agents operating in distinct
domains to benefit from each other's learned context by design.

The empirical evidence supports the hybrid pipeline as a minimum viable requirement
(nDCG@10 $0.5213 \pm 0.0004$ vs $0.0123 \pm 0.0000$ for FTS-only on natural-language
queries, n=50 3-run mean; absolute gap 50.9~pp). The pre-registered hypotheses for pain
weighting, cross-encoder comparison, and cross-corpus generalization are under evaluation
in sprint W2--W3; results will be published in the arXiv preprint at submission. Note that
the current nDCG@10 of $0.5213 < 0.6$, which keeps the D01 cross-encoder reranker gated
per \S6.5 until the threshold is met in future work.

Beyond nox-mem specifically, pain-weighted salience and shadow discipline are
\textbf{transferable concepts}. Any persistent memory system---regardless of implementation
stack---can adopt a severity annotation field and enforce a shadow validation gate before
ranking changes activate. These ideas require no new model, no new architecture, and no
GPU. They require only the discipline to instrument what already exists and the patience to
watch before activating.

The incidents are in the log. The log is in the schema. The schema is in this paper.

% ---------------------------------------------------------------
% BIBLIOGRAPHY
% ---------------------------------------------------------------
\bibliographystyle{unsrt}
\bibliography{../refs}

% ---------------------------------------------------------------
% APPENDICES
% ---------------------------------------------------------------
\appendix

\section{Implementation Details}

\subsection{TypeScript Stack}

The system is implemented in TypeScript (Node.js 22, strict mode, ESM modules) with the
following primary dependencies:

\begin{itemize}
  \item \texttt{better-sqlite3} --- synchronous SQLite interface; all DB operations are
    single-file transactions
  \item \texttt{sqlite-vec} --- vector extension for SQLite enabling cosine similarity
    search over 3072-dimensional Gemini embeddings
  \item \texttt{@google/generative-ai} --- Gemini API client for both embedding
    (Gemini embedding-001, 3072d) and LLM extraction (Gemini 2.5 Flash Lite for agent
    infra; Gemini 2.5 Flash for KG extraction)
  \item \texttt{inotifywait} --- filesystem watch for automatic ingest on file change
\end{itemize}

\textbf{Schema migration history (v1--v12).}

\begin{table}[h]
  \centering
  \begin{tabular}{ll}
    \toprule
    Version & Key change \\
    \midrule
    v1--v3 & Initial chunks + FTS5 design \\
    v4--v5 & \texttt{vec\_chunks} + \texttt{vec\_chunk\_map} (sqlite-vec) \\
    v6--v7 & \texttt{kg\_entities} + \texttt{kg\_relations} \\
    v8     & \texttt{retention\_days} typed retention policy \\
    v9     & \texttt{pain} field (REAL DEFAULT 0.2) \\
    v10    & \texttt{section} + \texttt{section\_boost} for entity file format \\
    v11    & \texttt{search\_telemetry} eval harness (+4 columns: \texttt{query\_text},
             \texttt{golden\_id}, \texttt{top\_chunk\_ids}, \texttt{top\_scores}) \\
    v12    & \texttt{ops\_audit} table + status enum enforcement triggers \\
    \bottomrule
  \end{tabular}
\end{table}

All migrations were applied to the production database without downtime.

\subsection{Entry Points}

\begin{itemize}
  \item \textbf{CLI:} \texttt{dist/index.js} (26+ subcommands including \texttt{search},
    \texttt{ingest}, \texttt{ingest-entity}, \texttt{reindex}, \texttt{vectorize},
    \texttt{kg-extract}, \texttt{reflect}, \texttt{crystallize})
  \item \textbf{MCP Server:} 16 tools including \texttt{nox\_mem\_search},
    \texttt{kg\_build}, \texttt{cross\_search}, \texttt{reflect}
  \item \textbf{HTTP API:} port 18802, endpoints
    \texttt{/api/\{health,search,kg,kg/path,agents,cross-kg,reflect,procedures\}}
    + \texttt{POST /api/crystallize}
\end{itemize}

% ------------------------------------------------------------------
\section{Shadow Case Study --- Fase~1.7b-b Salience Activation}

\subsection{Timeline}

\begin{itemize}
  \item \textbf{2026-04-18:} Pain-weighted salience formula implemented;
    \texttt{NOX\_SALIENCE\_MODE=shadow} set in production
  \item \textbf{2026-04-18 to 2026-04-25:} Seven-day shadow telemetry collection
  \item \textbf{2026-04-25:} Distribution analysis: 191 promotion candidates, 16,608
    review candidates, 45,743 archive candidates; distribution shows clear separation
    across tiers
  \item \textbf{2026-04-25:} Decision to advance to \texttt{NOX\_SALIENCE\_MODE=active};
    activation logged in \texttt{ops\_audit}
\end{itemize}

\subsection{Telemetry Summary}

\begin{figure}[h]
  \centering
  \includegraphics[width=0.85\linewidth]{figures/figure2-salience-pipeline.pdf}
  \caption{Salience pipeline. Each chunk is annotated at ingest with three
    signals---\texttt{retention\_days}, \texttt{section}, and \textbf{pain}
    (0.1 trivial to 1.0 prod-outage). Salience = $\texttt{recency} \times \texttt{pain}
    \times \texttt{importance}$ is computed and exposed via \texttt{/api/health.salience}
    for at least seven days before activation.}
  \label{fig:salience}
\end{figure}

\begin{table}[h]
  \caption{Telemetry distribution during Fase~1.7b-b shadow period.}
  \centering
  \begin{tabular}{lll}
    \toprule
    Tier & Count & \% of total \\
    \midrule
    Promoted (\texttt{new\_score} $\geq$ threshold) & 191 & 0.30\% \\
    Review (threshold range) & 16,608 & 25.88\% \\
    Archive (\texttt{new\_score} $<$ threshold) & 45,743 & 71.26\% \\
    \textbf{Total shadow observations} & \textbf{64,180+} & --- \\
    \bottomrule
  \end{tabular}
  \begin{minipage}{0.9\linewidth}
    \vspace{4pt}
    \footnotesize
    The 3-tier distribution reflects the expected long-tail structure of an operational
    corpus---most chunks are low-salience documentation, a minority are high-salience
    incident lessons.
  \end{minipage}
\end{table}

\subsection{Counterfactual: Incident 2026-04-25}

The incident of 2026-04-25 occurred the same day the shadow period ended---a coincidence
that underscores the motivation for the shadow gate. The reindex operation that damaged 183
entity records is precisely the class of ranking-affecting change that shadow mode is
designed to catch: it altered \texttt{section} and \texttt{section\_boost} annotations,
which feed directly into the salience computation. Had the post-incident reindex run before
the shadow gate was in place, the damage would have been invisible until a user noticed
degraded retrieval quality for entity queries.

\subsection{Activation Decision Rationale}

The distribution analysis showed that the pain-weighted salience formula was doing the
intended work: high-pain chunks (incident lessons, security decisions, production outage
records) consistently landed in the promotion tier, while trivial notes and meeting
summaries landed in the archive tier. The formula was activated because the distribution
matched the design intent, not because the absolute nDCG numbers improved (those were
measured separately in R01b).

% ------------------------------------------------------------------
\section{Reviewer-Friendly Feature Comparison Table}

This table summarizes the architectural features of nox-mem against the seven most closely
related systems. The scoring rubric follows the five dimensions identified in \S1.2: (1)
native knowledge graph with typed edges, (2) hybrid retrieval combining lexical and
semantic layers, (3) published evaluation harness with standard IR metrics, (4) multi-agent
shared context, (5) shadow-validated ranking discipline.

\begin{table}[h]
  \caption{Detailed architectural feature comparison (reviewer-friendly version of
    Table~\ref{tab:comparison}).}
  \centering
  \small
  \begin{tabular}{lllllll}
    \toprule
    System & KG native & Hybrid retr. & Eval harness & Multi-agent & Shadow disc. & Score \\
    \midrule
    \textbf{nox-mem (this work)} & Yes --- closed-enum & Yes --- FTS5+Gemini+RRF & Yes --- nDCG/MRR/Recall & Yes --- shared canonical & Yes --- $\geq$7d & \textbf{5/5} \\
    GraphRAG \cite{edge2024graphrag} & Yes + comm.\ det. & Partial --- via KG & No & No & No & 1.5/5 \\
    MemGPT/Letta \cite{packer2023memgpt} & No & Partial --- embed. & No & Yes --- per-agent & No & 1.5/5 \\
    Mem0 \cite{chhikara2025mem0} & Optional (v2) & No --- vector-only & Partial --- LOCOMO & Partial --- user\_id & No & 1.5/5 \\
    A-MEM \cite{xu2025amem} & Partial --- ZK & Partial --- semantic & No & No & No & 1.0/5 \\
    HiRAG \cite{huang2025hirag} & Yes --- hierarchical & Yes --- multi-level & Partial & No & No & 2.5/5 \\
    Cognee \cite{topoteretes2024cognee} & Yes --- ECL & Yes --- hybrid & Partial & Optional & No & 3.0/5 \\
    LangChain Memory & No & No --- key-value & No & Partial --- session & No & 0.5/5 \\
    \bottomrule
  \end{tabular}
\end{table}

% ------------------------------------------------------------------
\section{Incident Case Study --- Formal Write-Up (E13)}

\subsection{Incident Summary}

\textbf{Date:} 2026-04-25, 22:03 BRT \\
\textbf{Severity:} HIGH (data integrity; no user-facing downtime) \\
\textbf{Root cause:} End-of-day cron job executed \texttt{nox-mem reindex} without
dry-run or pre-operation snapshot. The generic \texttt{ingestFile()} path processed 183
entity files, stripping \texttt{section}, \texttt{retention\_days}, and
\texttt{section\_boost} annotations. \\
\textbf{Recovery:} Manual re-ingestion via \texttt{ingest-entity} for all 183 files;
\texttt{withOpAudit()} wrapper retroactively applied as preventive measure. \\
\textbf{Time to detect:} $\approx$18 minutes (via \texttt{/api/health.sectionDistribution}
check). \\
\textbf{Time to recover:} $\approx$45 minutes.

\subsection{Contributing Factors}

\begin{enumerate}
  \item The end-of-day cron script (cron ID \texttt{ee15b430}, 22:00 BRT, step 11)
    invoked \texttt{nox-mem reindex} without the \texttt{-{}-dry-run} flag.
  \item The reindex command used the generic \texttt{ingestFile()} router, which does not
    distinguish entity files from plain markdown.
  \item No pre-operation snapshot was taken; the daily backup (02:00 BRT) had not yet run
    for the day.
  \item No alert was configured for \texttt{section} annotation coverage drop.
\end{enumerate}

\subsection{Changes Implemented}

\begin{table}[h]
  \centering
  \begin{tabular}{lll}
    \toprule
    Change & Code artifact & Description \\
    \midrule
    F02 op-audit & \texttt{src/lib/op-audit.ts} & \texttt{withOpAudit()} wrapper: VACUUM INTO atomic snapshot \\
    A2 ingest router & \texttt{src/lib/ingest-router.ts} & \texttt{routeIngest()} dispatches entity files correctly \\
    A5 dry-run & \texttt{reindex.ts}, \texttt{consolidate.ts} & \texttt{-{}-dry-run} produces JSON preview without DB mutation \\
    Schema invariant canary & \texttt{check-schema-invariants.sh} & Cron \texttt{*/15min}; Discord alert on deviation \\
    Cron patch & \texttt{ee15b430} step 11 & Replaced \texttt{reindex} with \texttt{consolidate} (entity-aware) \\
    \bottomrule
  \end{tabular}
\end{table}

\subsection{Lessons Formalized in Architecture}

The \texttt{withOpAudit()} module encodes the lesson that destructive operations must
create a point-in-time snapshot before execution and log their outcome to an append-only
audit table (\texttt{ops\_audit}). The audit table's \texttt{status} field is validated
by DB triggers against a closed enum (\texttt{started}, \texttt{success},
\texttt{failed}, \texttt{crashed}); the triggers block DELETE and UPDATE on rows with
terminal status. Recovery via \texttt{safeRestore()} validates \texttt{user\_version}
match before restoring---a safeguard motivated by a separate incident where a stale WAL
file caused silent corruption on a naive \texttt{cp} restore.

This incident, and the five others documented in \texttt{docs/INCIDENTS.md}, are the
operational substrate from which the paper's architectural contributions were extracted.
They are included here not to confess failure but to demonstrate that the contributions
are grounded in production evidence rather than synthetic benchmark design.

\fi % end \iffalse — monolith §4-appendices suppressed

\end{document}
